\documentclass[a4paper,
fontsize=11pt,
%headings=small,
oneside,
numbers=noperiodatend,
parskip=half-,
bibliography=totoc,
final
]{scrartcl}

\usepackage[babel]{csquotes}
\usepackage{synttree}
\usepackage{graphicx}
\setkeys{Gin}{width=.4\textwidth} %default pics size

\graphicspath{{./plots/}}
\usepackage[ngerman]{babel}
\usepackage[T1]{fontenc}
%\usepackage{amsmath}
\usepackage[utf8x]{inputenc}
\usepackage [hyphens]{url}
\usepackage{booktabs} 
\usepackage[left=2.4cm,right=2.4cm,top=2.3cm,bottom=2cm,includeheadfoot]{geometry}
\usepackage[labelformat=empty]{caption} % option 'labelformat=empty]' to surpress adding "Abbildung 1:" or "Figure 1" before each caption / use parameter '\captionsetup{labelformat=empty}' instead to change this for just one caption
\usepackage{eurosym}
\usepackage{multirow}
\usepackage[ngerman]{varioref}
\setcapindent{1em}
\renewcommand{\labelitemi}{--}
\usepackage{paralist}
\usepackage{pdfpages}
\usepackage{lscape}
\usepackage{float}
\usepackage{acronym}
\usepackage{eurosym}
\usepackage{longtable,lscape}
\usepackage{mathpazo}
\usepackage[normalem]{ulem} %emphasize weiterhin kursiv
\usepackage[flushmargin,ragged]{footmisc} % left align footnote
\usepackage{ccicons} 
\setcapindent{0pt} % no indentation in captions
\usepackage{xurl} % Breaks URLs

%%%% fancy LIBREAS URL color 
\usepackage{xcolor}
\definecolor{libreas}{RGB}{112,0,0}

\usepackage{listings}

\urlstyle{same}  % don't use monospace font for urls

\usepackage[fleqn]{amsmath}

%adjust fontsize for part

\usepackage{sectsty}
\partfont{\large}

%Das BibTeX-Zeichen mit \BibTeX setzen:
\def\symbol#1{\char #1\relax}
\def\bsl{{\tt\symbol{'134}}}
\def\BibTeX{{\rm B\kern-.05em{\sc i\kern-.025em b}\kern-.08em
    T\kern-.1667em\lower.7ex\hbox{E}\kern-.125emX}}

\usepackage{fancyhdr}
\fancyhf{}
\pagestyle{fancyplain}
\fancyhead[R]{\thepage}

% make sure bookmarks are created eventough sections are not numbered!
% uncommend if sections are numbered (bookmarks created by default)
\makeatletter
\renewcommand\@seccntformat[1]{}
\makeatother

% typo setup
\clubpenalty = 10000
\widowpenalty = 10000
\displaywidowpenalty = 10000

\usepackage[colorlinks, linkcolor=black,citecolor=black, urlcolor=libreas,
breaklinks= true,bookmarks=true,bookmarksopen=true]{hyperref}
\usepackage{hyperxmp}
\usepackage{breakurl}

%meta

%meta

\fancyhead[L]{Redaktion LIBREAS\\ %author
LIBREAS. Library Ideas, 48 (2026). % journal, issue, volume.
\href{https://doi.org/10.18452/37907}{\color{black}https://doi.org/10.18452/37907}
{}} % doi 
\fancyhead[R]{\thepage} %page number
\fancyfoot[L] {\ccLogo \ccAttribution\ \href{https://creativecommons.org/licenses/by/4.0/}{\color{black}Creative Commons BY 4.0}}  %licence
\fancyfoot[R] {ISSN: 1860-7950}

\title{\LARGE{Das liest die LIBREAS, Nummer \#17 (Frühjahr 2026)}}% title
\author{Redaktion LIBREAS} % author

\setcounter{page}{1}

\hypersetup{%
      pdftitle={Das liest die LIBREAS, Nummer \#17 (Frühjahr 2026)},
      pdfauthor={Redaktion LIBREAS},
      pdfcopyright={CC BY 4.0 International},
      pdfsubject={LIBREAS. Library Ideas, 48 (2026)},
      pdfkeywords={Literaturübersicht, Bibliothekswissenschaft, Informationswissenschaft, Bibliothekswesen, Rezension, literature overview, library science, information science, library sector, review},
      pdflicenseurl={https://creativecommons.org/licenses/by/4.0/},
      pdfcontacturl={http://libreas.eu},
      baseurl={},
      pdflang={de},
      pdfmetalang={de},
      pdfdoi={10.18452/37907},
      pdfurl={https://doi.org/10.18452/37907}
     }



\date{}
\begin{document}

\maketitle
\thispagestyle{fancyplain} 

%abstracts

%body
Beiträge von Karsten Schuldt (ks), Viola Voß (vv), Eva Bunge (eb)

\section{1. Zur Kolumne}\label{zur-kolumne}

Ziel dieser Kolumne ist es, eine Übersicht über die in der letzten Zeit
erschienene bibliothekarische, informations- und
bibliothekswissenschaftliche sowie für diesen Bereich interessante
Literatur zu geben. Enthalten sind Beiträge, die der LIBREAS-Redaktion
oder anderen Beitragenden als relevant erschienen.

Themenvielfalt sowie ein Nebeneinander von wissenschaftlichen und
nicht-wissenschaftlichen Ansätzen wird angestrebt und auch in der Form
sollen traditionelle Publikationen ebenso erwähnt werden wie
Blogbeiträge oder Videos beziehungsweise TV-Beiträge.

Gerne gesehen sind Hinweise auf erschienene Literatur oder Beiträge in
anderen Formaten. Diese bitte an die Redaktion richten. (Siehe
\href{http://libreas.eu/about/}{Impressum}, Mailkontakt für diese
Kolumne ist
\href{mailto:zeitschriftenschau@libreas.eu}{\nolinkurl{zeitschriftenschau@libreas.eu}}.)
Die Koordination der Kolumne liegt bei Karsten Schuldt, verantwortlich
für die Inhalte sind die jeweiligen Beitragenden. Die Kolumne
unterstützt den Vereinszweck des LIBREAS-Vereins zur Förderung der
bibliotheks- und informationswissenschaftlichen Kommunikation.

LIBREAS liest gern und viel Open-Access-Veröffentlichungen. Wenn sich
Beiträge dennoch hinter einer Bezahlschranke verbergen, werden diese
durch \enquote{{[}Paywall{]}} gekennzeichnet. Zwar \linebreak macht das Plugin
\href{http://unpaywall.org/}{Unpaywall} das Finden von legalen
Open-Access-Versionen sehr viel einfacher. Als Service an der
Leserschaft verlinken wir jedoch auch direkt OA-Versionen, die wir vorab
finden konnten. Für alle Beiträge, die dann immer noch nicht frei
zugänglich sind, empfiehlt die Redaktion (neben
\href{http://unpaywall.org/}{Unpaywall}) die Browser-Plugins
\href{https://openaccessbutton.org/}{Open Access Button} oder
\href{https://core.ac.uk/services/discovery/}{CORE} zu nutzen sowie auf
dem favorisierten Social-Media-Kanal mit
\href{https://mastodon.social/tags/icanhazpdf}{\#icanhazpdf}, um Hilfe
bei der legalen Dokumentenbeschaffung zu bitten.

Die bibliographischen Daten der besprochenen Beiträge aller Ausgaben
dieser Kolumne finden sich in der öffentlich zugänglichen Zotero-Gruppe:
\url{https://www.zotero.org/groups/4620604/libreas_dldl/library}.

\section{2. Artikel und
Zeitschriftenausgaben}\label{artikel-und-zeitschriftenausgaben}

\section{2.1 Vermischte Themen}\label{vermischte-themen}

\emph{Bibliothèque(s) : revue de l\textquotesingle Association des
bibliothécaires français (ABF)}. (2025) 104

Im Jahr 2021 wurde die französische Zeitschrift \emph{Bibliothèque(s)}
eingestellt. Sie war bis dahin die wichtigste praxisorientierte
Zeitschrift für Öffentliche Bibliotheken in Frankreich gewesen.
Vergleichbar war sie mit der \emph{Forum Bibliothek und Information
(BuB)}, den \emph{Büchereiperspektiven} und mehr noch -- weil sie ebenso
eigentlich für alle Bibliotheken gedacht ist, sich aber vor allem auf
die Öffentlichen fokussierte -- mit der \emph{bibliosuisse info}.
Herausgegeben wurde sie vom Verband der französischen
Bibliothekar*innen.

Angesichts dessen, dass in diesem Zeitraum noch mehr französische
bibliothekarische Publikationen eingestellt wurden, schien dies einen
problematischen Trend anzuzeigen. Insoweit ist es positiv hervorzuheben,
dass die \emph{Bibliothèque(s)} seit Anfang 2025 wieder regelmässig
erscheint. Das Format wurde angepasst, jetzt rund 60 Seiten auf A4.
Erscheinen soll sie zweimal im Jahr.

Im Editorial der ersten Ausgabe (Mai 2025) wird zwar geschrieben, dass
es ein neues Konzept gäbe, aber wie dieses aussieht, wird nicht gesagt.
Inhaltlich ist diese Ausgabe gegliedert in ein Dossier -- also eine
Anzahl von Artikeln und Interviews zu einem Thema, hier die Evaluation
von Bibliotheken, wobei es sowohl um Statistik, um das, was in Deutsch
Gemeinwohlanalyse genannt würde, um die Interpretation von Daten in
Bibliotheken, als auch um eine Enquête über alle Bibliotheken
Frankreichs geht --, eine Reihe anderer Artikel, der Rubrik
Bibliotourisme, inklusiver einer Bilderstrecke -- konkret über eine
Bibliothek in Barcelona -- sowie Meldungen über die Arbeit des Verbandes
und seiner Kommissionen.

Zudem werden die Präsidentin der International Federation of Library
Associations and Institutions (IFLA) vorgestellt und drei Rezensionen
publiziert. Das alles ist, für den französischen Kontext, wo Texte und
Bücher oft sehr lang sind, äusserst kurz gehalten. Die meisten Texte
sind zwei bis drei Seiten lang, die Rezensionen nur je eine halbe Seite.
Zudem ist vieles sehr \enquote{hands-on}, also mit einem Blick daraufhin
geschrieben, dass es direkt in Bibliotheken umgesetzt werden kann. Zu
vermuten ist, dass dieser Aufbau in den nächsten Ausgaben beibehalten
wird. So französische Sprachkenntnisse vorhanden sind, sollte die
\emph{Bibliotheque(s)} jetzt (wieder) in die fachliche
Zeitschriftenschau von Bibliothekar*innen im Öffentlichen
Bibliothekswesen aufgenommen werden. (ks)

\begin{center}\rule{0.5\linewidth}{0.5pt}\end{center}

Golijanin, Aleksandar (2025). \emph{What Language Are We Speaking?:
Marketing Information Literacy on University Library Websites}. In:
Communications in Information Literacy 19 (2025) 1: 51--68,
\url{https://doi.org/10.15760/comminfolit.2025.19.1.4}

In dieser Studie wird gefragt, wie Bibliotheken an den grössten
kanadischen Universitätsbibliotheken ihre Angebote im Bereich
Information Literacy gegenüber Forschenden und Lernenden beschreiben. Es
ist offenbar -- obgleich es in Kanada und den USA eine Tradition von
Beschäftigung mit dem Thema gibt, die noch umfassender und älter ist,
als im DACH-Raum -- in der Praxis schwierig, dem
nicht-bibliothekarischen Personal an den Hochschulen klar zu machen,
welche Angebote (Unterricht, Einführungen, Beratungen und so weiter) die
Bibliotheken machen. Hier wird die Sprache verglichen, welche die
Bibliotheken dafür benutzen. Dabei zeigt sich, dass ein Grossteil doch
den Begriff \enquote{Information Literacy} nutzt, der aber gleichzeitig immer
wieder erklärt werden muss. Zudem gibt es drei alternative
Beschreibungen (\enquote{research skills}, \enquote{find, evaluate, and use information}
sowie \enquote{critical information use}), welche ebenso oft genutzt werden,
wohl -- so die Vermutung -- weil die Bibliotheken den Eindruck haben,
dass dies von den Lehrenden und Forschenden besser verstanden wird.

Zudem zeigt sich eine geographische Verteilung. In Ontario und Quebec --
wo sich tendenziell auch grössere und forschungsstärkere Universitäten
befinden -- wird der Begriff \enquote{Information Literacy} viel weniger genutzt
als in den anderen Landesteilen. Die Studie äussert keine Vermutung,
warum dies so ist, ruft aber zu einer Untersuchung dieses Phänomens auf.
(ks)

\begin{center}\rule{0.5\linewidth}{0.5pt}\end{center}

Hall, Suzanna (2025). \emph{\enquote*{Please dispose of properly}: the
context and continuance of fashion ephemera at the International Library
of Fashion Research}. In: Art Libraries Journal 50 (2025) 1: 45--51,
\url{https://doi.org/10.1017/alj.2025.10} {[}Paywall{]}

Die \enquote{International Library of Fashion Research} ist eine Einrichtung,
angesiedelt im norwegischen Nationalmuseum in Oslo, die sich dem
Sammeln, Erschliessen und Zugänglichmachen von ephemeren Materialien
widmet, die in der Modeindustrie produziert werden. Damit sind zum
Beispiel Einladungen zu Modeschauen, Entwürfe von Kleidung oder
saisonale Kataloge, die nur für einen Zeitraum von einigen Monaten
\enquote{gültig} sind, gemeint. Viele dieser Materialien sind, da sie aus der
Modeindustrie stammen und immer auch Aufmerksamkeit erzeugen sollen,
gesondert gestaltet oder ausgestattet. In Oslo wurde nun auf der Basis
einer privaten Sammlung, welche dem Museum gespendet wurde, die genannte
Bibliothek errichtet, welche seit Ende 2022 für die Öffentlichkeit im
Museumsgebäude zugänglich ist. Der Artikel stellt diese Einrichtung auf
der Basis von Interviews mit Beteiligten vor.

Es wird begründet, warum die Materialien gesammelt werden (nämlich als
Material für die Modegeschichte), zudem wird auf die Schwierigkeiten bei
der Erhaltung, die solche Materialien machen, eingegangen. Das ist alles
recht gut nachvollziehbar. Was im Artikel aber hervorsticht, ist, dass
sich hier \enquote{Museumsleute} dazu äussern, was sie sich unter einer
Bibliothek, einem Archiv und einem Museum vorstellen. Es wird betont,
dass es explizit eine Bibliothek sei, weil die Materialien zugänglich
wären. Laut den Stimmen, die im Artikel angeführt werden, wäre das bei
Archiven nicht der Fall (etwas, was selbstverständlich falsch ist --
auch bei Archiven besteht das Ziel ja letztlich in der Nutzung der
Sammlung). Gleichzeitig orientieren sie sich an dem Bild einer
Bibliothek, die vor allem bewahrt. So wird nicht erwähnt, ob die
Bibliothek überhaupt ein aktives Bestandsmanagement -- im Sinne des
Aufbaus von spezifischen Beständen und der Pflege inklusive des
Aussonderns und Ersetzens von Medien -- betreibt. Dabei würde sich die
Fachliteratur zur Modegeschichte für einen solchen Forschungsbestand,
neben der Sammlung, direkt anbieten. Vielmehr wird hier offenbar mit
\enquote{Bibliothek} eine sehr spezifische Einrichtungen gemeint, die sich sonst
an Büchern orientieren würde (was selbstverständlich gar nicht mehr
stimmt). Es ist auch nicht klar, ob beim Aufbau der Bibliothek überhaupt
bibliothekarisches Personal herangezogen oder aber anders auf
bibliothekarische Kompetenzen zurückgegriffen wurde. Eher wird der
Eindruck vermittelt, dass hier \enquote{Museumsleute} eine Einrichtung aufgebaut
haben, die vor allem \enquote{kein Archiv} und \enquote{kein Museum} sein soll -- und
deshalb \enquote{Bibliothek} genannt wird. (ks)

\begin{center}\rule{0.5\linewidth}{0.5pt}\end{center}

Braisher, Tamsin ; Fitchett, Deborah (2025). \emph{The New Zealand
Thesis Project: Connecting a Nation's Dissertations Using Wikidata}. In:
Journal of Librarianship and Scholarly Communications 13 (2025) 1,
\url{https://doi.org/10.31274/jlsc.18295}

Zwischen Ende 2021 und Ende 2022 wurden in einem Projekt, an dem alle
Universitäten und Hochschulinstitute des Landes beteiligt waren, die
Metadaten zu den Dissertationen, die in Aotearoa/Neuseeland seit den
1960er Jahren publiziert wurden (und zu denen Katalog- oder andere
Metadaten vorlagen) in die Wikipedia eingefügt. Einmal als gesamtes
Datenset für die Wikidata und dann noch als Daten auf den
Wikipedia-Seiten zu den Forschenden und Betreuenden, falls diese
existierten. Wenn die Dissertationen in Repositorien vorlagen -- also
wenn sie dort an sich publiziert wurden oder wenn Digitalisate
angefertigt worden waren --, wurden die URIs zu diesen in den Daten
ergänzt. Zudem wurden aus den Daten Visualisierungen erzeugt. Der
Artikel stellt einen ausführlichen Projektbericht dar.

Grundsätzlich wurden in dem Projekt alle Schritte begangen, die in einem
solchen wohl notwendig sind: Es wurden möglichst alle Einrichtungen des
Landes \enquote{an Bord geholt}, sich auf ein Metadatenschema geeinigt, Daten
wurden geliefert und gesäubert, dann in Wikidata eingefügt.
Anschliessend wurden die betreffenden Wikipedia-Seiten recherchiert und
ergänzt. Das alles ist im Artikel recht kleinteilig geschildert. Dabei
zeigt sich vor allem, wie wichtig die Planung und dann die Datenpflege
selber ist -- also die \enquote{Datenarbeit} --, um am Ende überhaupt
Visualisierungen generieren zu können. Durch diese Schilderung lässt
sich der Text sehr gut als Einstieg und Lehrmaterial nutzen --
beispielsweise, wenn am Anfang eines bibliotheks- und
informationswissenschaftlichen Studiums die Notwendigkeit dieser
\enquote{Datenarbeit} vermittelt werden soll. Er ist recht eingängig
geschrieben, schildert alle notwendigen Schritte, zeigt Ergebnisse, aber
auch Herausforderungen. (ks)

\begin{center}\rule{0.5\linewidth}{0.5pt}\end{center}

Hoffmann, Kristin (2025). \emph{Exploring Canadian academic librarians'
and archivists' research outside librarianship and archival studies}.
In: Partnership, The Canadian Journal of Library and Information
Practice and Research, 20 (2025) 1,
\url{https://doi.org/10.21083/partnership.v20i1.8141}

In Wissenschaftlichen Bibliotheken Kanadas gibt es oft -- aber nicht
überall -- die Erwartung an Bibliothekar*innen, zu forschen und zu
publizieren. Teilweise wird dies bei Bewertungen ihrer Arbeit
einbezogen, in vielen Fällen wird es auch direkt von den Einrichtungen
unterstützt. In einer Umfrage, die hier ausgewertet wird, geht Hoffmann
auf einen spezifischen Bereich ein, nämlich auf Forschung, die zwar von
Bibliothekar*innen (und Archivar*innen) aus Kanada betrieben wird, sich
aber nicht in der Bibliotheks- oder Archivwissenschaft verorten lässt.
Es wird gefragt, in welchen Forschungsdisziplinen dies passiert, warum
die Bibliothekar*innen und Archivar*innen dies tun und auch, ob es dafür
Barrieren gibt.

Im Ergebnis zeigt sich, dass dies alles ein weites, unerforschtes Feld
ist. Die Disziplinen, in denen geforscht und publiziert wird, sind
zahlreich; die Gründe vielfältig. Auch die institutionellen Strukturen
und Erwartungen sind im ganzen Land sehr unterschiedlich. Viele der
Antwortenden betonten, dass es für ihre Arbeit und ihre
Arbeitsplatzzufriedenheit wichtig wäre, auch ausserhalb der Bibliotheks-
und Archivwissenschaft zu forschen. Aber ein Hauptgrund scheint doch
immer wieder das persönliche Interesse zu sein.

Interessant ist, dass sich zeigt, dass mehr Männer als Frauen --
obgleich auch in Kanada das Bibliothekswesen \enquote{weiblich geprägt} ist --
so forschen. Sowohl Kommentare, die in der Umfrage erhoben wurden, als
auch Hoffmann in der Auswertung selber deuten darauf hin, dass dies auch
für andere Kriterien gilt: Obgleich es offenbar in Kanada recht viele
Möglichkeiten zur forschenden Tätigkeit gibt, wenn jemand in
Wissenschaftlichen Bibliotheken angestellt ist, sind es dann
insbesondere Personen aus nicht-marginalisierten Gruppen, welche diese
Möglichkeiten haben und wahrnehmen. (ks)

\begin{center}\rule{0.5\linewidth}{0.5pt}\end{center}

Jonchery, Anne ; Octobre, Sylvie (2025). \emph{Le numérique sauvera-t-il
la lecture chez les jeunes?.} In: InterCDI (2025) Septembre-Octobre:
5--12 {[}gedruckt{]}

Ortin, Manon (2025). \emph{Les jeunes et la lecture numérique: entre
diversité des formats et nouveaux modes de consommation}. In: InterCDI
(2025) Septembre-Octobre: 14--22 {[}gedruckt{]}

Beide Artikel erschienen im Rahmen eines Schwerpunkts zum Thema
\enquote{Lesen} der Zeitschrift für französische Schulbibliotheken.
Insoweit geht es immer um das Lesen von schulpflichtigen Kindern und
Jugendlichen, so wie es sich in den letzten Jahren in Frankreich
entwickelt hat. Berichtet wird in beiden Artikeln darüber, wie diese
Kinder und Jugendlichen sich dem Lesen von digitalen Medien zugewandt
haben. Die Daten dazu stammen aus mehreren Umfragen und
Beobachtungsreihen zum Lesen, die in Frankreich staatlicherseits erhoben
werden.

Beide Artikel machen zwei Hauptaussagen: Erstens hat sich das Lesen von
digitalen Medien (oder der Konsum von Audiobooks) bei den Kindern und
Jugendlichen als Alltagspraxis durchgesetzt, aber nicht vollständig.
Zweitens ist dies nicht nur ständig in Entwicklung, sondern zeigt auch
eine grosse Bandbreite an Medienformen, Inhalten und Nutzungsweisen. Es
ist nicht möglich, einen klaren Trend zu zeigen, ausser den, dass sich
diese Diversität immer weiter ausdehnt. Es ist nicht wirklich möglich,
von einer gemeinsamen \enquote{Medienwelt} der Kinder und Jugendlichen
zu sprechen. Was zunimmt, ist die aktive Partizipation von Kindern und
Jugendlichen, indem sie in Communities Kommentare verfassen, ihre
Meinungen per Social Media teilen oder auch selber Literatur
veröffentlichen. Sichtbar ist allerdings für Leser*innen aus dem
DACH-Raum auch der kulturelle Unterschied: In Frankreich sind Comics,
Webtoons und Manga heute massiv prägende Medienformen, deren Wirkung und
Verbreitung sich nicht mit der im DACH-Raum vergleichen lässt. Sie
prägen auch den Medienkonsum in Frankreich (oder in der Schweiz der
Romandie) viel mehr.

Die Daten und wohl auch Trends im DACH-Raum werden wohl andere sein. Die
Artikel machen deshalb auch schmerzlich sichtbar, dass regelmässige
Datensammlungen zum Lektüreverhalten diesen Umfangs, die nicht aus dem
Umfeld von Verlagen selber stammen, im DACH-Raum praktisch fehlen. (ks)

\subsection{2.2 Open Access}\label{open-access}

Jordan, Jennifer ; Solon, Blair ; Beene, Stephanie (2025).
\emph{Identifying Open Access Practices in Librarianship Journals}. In:
Journal of Librarianship and Scholarly Communications 13 (2025),
\url{https://doi.org/10.31274/jlsc.17778}

Jordan, Jennifer ; Solon, Blair ; Beene, Stephanie (2025). \emph{Open
access practices of selected library science journals}. {[}Dataset{]}
Dryad, \url{https://doi.org/10.5061/dryad.pvmcvdnt3}

In dieser Studie wurde untersucht, wie \enquote{offen} (im Sinne der
FAIR-Prinzipien) aktuell die Zeitschriften des Bibliothekswesens sind --
zumindest, was die wichtige Einschränkung hier ist, wenn sie
Informationen und Abstracts in Englisch publizieren. Die Frage ist
selbstverständlich relevant, da die Differenz zwischen dem Anspruch von
Bibliotheken, Open Access sowie andere Formen von Open Science zu
fördern und der gleichzeitigen Existenz von Closed-Access-Zeitschriften
im Bibliothekswesen immer wieder auffällt. Die Autor*innen der Studie
erstellten eine Liste von bibliothekswissenschaftlichen und
bibliothekarischen Zeitschriften, die entweder in der Datenbank LISA
oder im DOAJ verzeichnet waren, einen Peer-Review-Prozess haben und dies
in Englisch erwähnen. Zudem ergänzten sie diese um weitere Titel, die
ihnen als Bibliothekar*innen selber bekannt waren, die aber weder in
LISA noch im DOAJ verzeichnet waren. Am Ende waren das 73 Zeitschriften
aus LISA, 61 aus dem DOAJ (mit 12, die in beiden Datenbanken verzeichnet
waren) und 11 ergänzte Titel. Anschliessend sammelten sie Informationen
über die herausgebenden Körperschaften, die \enquote{OA-Farbe}, Angaben zu Open
Access und Open Data Policies sowie zu eventuell anfallenden Gebühren
(also APCs).

In der Auswertung der Daten zeigt sich, dass diese Zeitschriften viel
öfter als nicht Open-Access-Modellen folgen und dabei vor allem dem
Diamond Open Access (also auch ohne anfallende Gebühren). Sie werden oft
von kommerziellen Verlagen publiziert, aber leicht öfter von Verbänden,
Hochschulen, Bibliotheken, Archiven oder \enquote{unabhängig} (was heisst von
Organisationen wie den LIBREAS. Verein). Es gibt Unterschiede zwischen
den Zeitschriften, die in LISA verzeichnet sind (eher kommerzielle
Verlage) und denen im DOAJ (eher unabhängig oder von nicht-kommerziellen
Körperschaften). Schwieriger war es offenbar, zu bestimmen, welche
Formen von Peer Review (also vor allem, ob es ein Open Peer Review gibt)
genutzt werden, ob es Open Data Policies gibt (hier zeigte sich, dass
viele grosse Verlage praktisch verlagsweit eine Policy durchsetzen) oder
welcher Höhe eventuell anfallende Gebühren sind. Alles in allem zeigt
sich aber ein klarer Trend zu offenen Publikationen im
Bibliotheksbereich. (ks)

\begin{center}\rule{0.5\linewidth}{0.5pt}\end{center}

Hopkins, Kira ; Sanders, Kevin (2025). \emph{Open access, open
infrastructures, and their funding: Learning from histories to more
effectively enhance diamond OA ecologies for books}. In: Journal of
Librarianship and Scholarly Communications 12 (2024) 2,
\url{https://doi.org/10.31274/jlsc.18284}

Dieser Artikel bietet eine Rundumsicht zu den Finanzierungsmodellen für
Open-Access-Bücher, die sich in den letzten Jahren etabliert haben,
wobei der Fokus auf Modellen liegt, die gerade keine gesonderten
Gebühren auf Seiten der Autor*innen erheben. Wer sich auf dem Feld
einigermassen auskennt, wird keine Neuigkeiten erfahren. Vielmehr ist
der Text als Einstieg in das Thema und die Modelle geeignet. (ks)

\begin{center}\rule{0.5\linewidth}{0.5pt}\end{center}

Rawlins, Ben ; Scott, Mitchell (2025). \emph{Open Access and Citation
Impact: Modality, Funding, Publisher, and Disciplinary Trends at the
University of Kentucky}. In: Library Resources \& Technical Services 69
(2025) 3, \url{https://doi.org/10.5860/lrts.69n3.8496}

Der \enquote{Open-Access-Zitationsvorteil} ist ein recht etabliertes Phänomen,
welches schon früh in der Forschung zu Open Access untersucht und immer
wieder nachgewiesen wurde: Im Vergleich zu wissenschaftlichen Artikeln,
die Closed Access publiziert werden, werden Open-Ac\-cess-Publikationen
mehr und früher zitiert. Abgeleitet wird daraus auch oft, dass sie wohl
mehr und früher gelesen werden, also in ihrer Funktion als
wissenschaftliche Publikation erfolgreicher sind. Die vorliegende Studie
untersucht dieses Phänomen anhand der Beiträge, die von Forschenden
einer spezifischen Hochschule -- der University of Kentucky --
veröffentlicht wurden und das auch in einem spezifischen
Hochschulkontext, dem US-amerikanischen, in welchem beispielsweise keine
Tradition von Konsortiallizenzen für digitale Medien auf nationaler
Ebene existieren, wie das in Europa der Fall ist. Die Daten lassen sich
also nicht einfach übertragen.

Im Ergebnis zeigt die Studie aber erstens, wie eine solche Auswertung
konkret stattfinden kann (also welche Methodik, welche Fragen und welche
Formen von Präsentation der Daten und Ergebnisse gewählt werden können)
und zweitens, dass der \enquote{Open-Access-Zitationsvorteil} differenzierter
ist. Während er grundsätzlich auch für die Publikationen aus der
University of Kentucky gilt, zeigt sich doch, dass er sehr
unterschiedlich ausgeprägt ist in den einzelnen Wissenschaftsdisziplinen
und auch Verlagen. Teilweise ist er gar nicht vorhanden, teilweise
werden Open-Access-Publikationen in \enquote{hybriden} Zeitschriften, die sowohl
Open als auch Closed Access veröffentlichen, mehr und schneller zitiert
als solche, die in \enquote{reinen} Open Access Journals veröffentlicht werden.
Aber teilweise ist er auch überdeutlich zu sehen. (ks)

\subsection{2.3 Bestandsmanagement}\label{bestandsmanagement}

Steele, Gemma ; Webster, Hayley (2025). \emph{Getting Inked? A Survey of
Current Institutional Marking Practices in Rare Books and Special
Collections}. In: RBM: A Journal of Rare Books, Manuscripts, and
Cultural Heritage 26 (2025) 1: 52--66,
\url{https://doi.org/10.5860/rbm.26.1.52}

Der Text berichtet von einer Umfrage unter Bibliotheken -- wohl, da sie
über betreffende Mailinglisten verbreitet wurde, vor allem in
englischsprachigen Ländern (aus Sicherheitsgründen wurden aber keine
geographischen Daten abgefragt) -- darüber, ob Bibliotheken ihre
Rara-Bestände -- also die seltenen Drucke und Handschriften -- mit einem
Bibliotheksstempel versehen oder nicht. Das klingt wie ein
nebensächliches Thema, es ist aber eine -- wie der Artikel zeigt -- seit
Jahrzehnten diskutierte Frage. Argumente gegen und für das Stempeln
werden offenbar regelmässig ausgetauscht: Dafür spricht, dass diese
Stempel vor allem als Diebstahlschutz dienen. Wenn sie vorhanden sind,
insbesondere so angebracht, dass sie kaum entfernt werden können, sei es
schwieriger, diese Objekte nach einem Diebstahl zu verkaufen. Dagegen
spricht, dass solche Stempel immer das Objekt selber verändern, also
einen Eingriff darstellen. Es gibt, wie im Text auch zu lernen ist, eine
Reihe von anderen Ansätzen zur Sicherung von Rara (unsichtbare Tinte,
Katalogisierung und Dokumentation von relevanten Eigenheiten der Rara
und anderes mehr) sowie einige weitere Argumente. Zudem wurden schon
mehrere übergreifende Richtlinien erlassen, die einerseits das Stempeln
empfehlen, aber es jeder Bibliothek selber überlassen, sich dafür oder
dagegen zu entscheiden.

Die Umfrage selber brachte jetzt, trotz 93 Antworten, wenig Neues. Die
Praxis ist offenbar in den verschiedenen Bibliotheken sehr
unterschiedlich, die Argumente sind kaum anders, als die in der
Literatur der letzten Jahrzehnte diskutierten. (Zu kritisieren ist aber,
dass die Literatur sich praktisch auf englischsprachige beschränkt.) Im
Ganzen liefert der Artikel aber einen interessanten Einblick in das
Thema. (ks)

\begin{center}\rule{0.5\linewidth}{0.5pt}\end{center}

Ho, Jessica (2025). \emph{There are Imposters Among Us -- Investigating
AI-Generated Picture Books in the Library}. In: Canadian School
Libraries Journal 32 (2025) 2,
\url{https://journal.canadianschoollibraries.ca/there-are-imposters-among-us-investigating-ai-generated-picture-books-in-the-library/}
{[}Re-Published{]}

In diesem kurzen Text gibt Ho, angestellt im Schulbibliothekssystem von
Toronto, eine Übersicht ihrer Gedanken zur steigenden Anzahl von
Kinderbüchern, die offenbar ohne Qualitätskontrolle unter Einsatz von
AI-Tools erstellt und über Anbieter wie Amazon verbreitet werden. Ho
gibt Hinweise, wie diese Bücher gestaltet sind und sich erkennen lassen.
Zudem wird besprochen, dass sich diese Bücher auf den ersten Blick oft
nicht von Selbst-Publikationen und Publikationen kleiner Verlage
unterscheiden, auf den zweiten Blick dann aber doch. Zuletzt verweist
der Text aber darauf, dass vor allem Bibliothekar*innen die
Verantwortung hätten \enquote{to make informed decisions and advocate for
intellectual property rights, valuable stories and diverse
representation}. Für die Erstellung von Kinderbüchern würden immer auch
die aktuell verfügbaren Tools genutzt werden und AI sei nur ein weiteres
Tool. (ks)

\begin{center}\rule{0.5\linewidth}{0.5pt}\end{center}

Wolf, Christa (2025). \emph{Musikinstrumente zum Ausleihen: Ein
Erfahrungsaustausch im Verbund der Öffentlichen Bibliotheken Berlins.}
In: Forum Musikbibliothek 46 (2025) 3: 6--10 {[}gedruckt{]}

Heß, Dina ; Mahn, Johanna (2025). \emph{Bibliothek der Dinge -- Mobile
Arbeitsmittel für Studierende in der Folkwang Bibliothek}. In: Forum
Musikbibliothek 46 (2025) 3: 10--16 {[}gedruckt{]}

Dargestellt wird in den beiden Artikeln vor allem, wie Musikinstrumente,
Stative und iPads, immer gedacht als Teil von \enquote{Bibliotheken der
Dinge}, in der Praxis in Bibliotheken gehandhabt werden. Der Artikel aus
Berlin geht auch auf die spezifischen Herausforderungen der
Vergabevorschriften ein (die in anderen Städten oder Ländern wieder
anders sind). Was sich bei beiden Texten zeigt, ist, dass Bibliotheken
Wege suchen und finden, um \enquote{Dinge} anzuschaffen, ausleihfertig
zu machen, zu katalogisieren und zu pflegen, aber dass das alles sehr
Stückwerk ist: Ständig wird von Bibliothekar*innen zur Eigeninitiative
gegriffen, auf die spezifische Interessen und Fähigkeiten des Personals
zurückgegriffen (oder Absprachen mit anderen getroffen, unter anderem
mit einer Instrumentenwerkstatt) und spezifische Lösungen gesucht, die
irgendwie in die bestehenden Systeme passen (beispielsweise die
Einführung einer \enquote{fiktiven Reihe} für die Katalogisierung,
\enquote{Musikinstrumente zum Ausleihen}, damit Instrumente überhaupt
als eigene Gruppe im Katalog gefunden werden können).

Bestandsarbeit ist offenbar, wenn es sich nicht um \enquote{normale
Medien} handelt, immer eine Art von Life-hack. (ks)

\begin{center}\rule{0.5\linewidth}{0.5pt}\end{center}

Mathews, Emilee ; Emerson, María Emerson (2025). \emph{From Niche to
Norm: A Case Study of Zines in a Circulating Collection.} In: Library
Resources \& Technical Services 69 (2025) 4,
\url{https://doi.org/10.5860/lrts.69n4.8550}

Zines, also von Einzelpersonen oder kleinen Gruppen oft per Hand
erstellte und in kleiner Auflage verbreitete Broschüren, stellen einen
eigenen, schwer zu fassenden Medientyp dar. Oft als Teil der
Kommunikation in einzelnen Subkulturen etabliert, folgen sie selten
expliziten Vorgaben, erscheinen oft unregelmässig oder auch nur einmal,
sind nirgends vollständig verzeichnet (was ihrem Zweck oft auch
zuwiderlaufen würde) und werden nur in Ausnahmefällen zum freien Verkauf
angeboten. Aber gleichzeitig, oder vielleicht sogar deshalb, sind sie
immer wieder von Interesse für die Forschung und auch für Bibliotheken.
Letzteres trifft offenbar noch mehr für US-amerikanische Bibliotheken zu
als für die im DACH-Raum. Dort stellen Zines-Bestände Sondersammlungen
dar, die teilweise regionale Bedeutung haben.

Der Artikel stellt eine dieser Sammlungen, die Library Social Justice
Zine Collection der University of Illinois Urbana-Champaign, aus
bibliothekarischer Sicht vor, gibt aber auch gleichzeitig einen
Überblick über die bisherigen Publikationen zum Thema sowie Ergebnisse
einer Umfrage. Das macht den Text sehr uneinheitlich -- eigentlich sind
es drei Artikel, die auch getrennt hätten veröffentlicht werden können.
Der Literaturteil und die Umfrage zeigt aber, dass sich \enquote{Zine
Sammlungen} in den US-amerikanischen Hochschulbibliotheken verbreitet
haben (ohne dafür Gründe angeben zu können) und genutzt werden. Erwähnt
wird auch die selbstorganisierte Gruppe hinter der Homepage
\url{https://www.zinelibraries.info/}. Gerade der Teil, welcher die
spezifische Sammlung vorstellt, macht aber auch klar, dass das
Charakteristikum von Zines, sich in Form, Inhalt und so weiter
kontinuierlich zu unterscheiden und zudem (oft) aus einfach Materialien
hergestellt zu sein, immer zu besonderen Lösungen für die
Katalogisierung, Aufbewahrung und Anbietung führt. Es ist viel Arbeit
für manchmal recht wenige Medien. (ks)

\begin{center}\rule{0.5\linewidth}{0.5pt}\end{center}

Duffus, Orolando (2025). \emph{Strategies for Assessing and Enhancing
Electronic Journal Collections: A Data-Driven Approach for Supporting
New Academic Programs and the College of Medicine}. In: International
Journals of Librarianship 10 (2025) 4, 5--16,
\url{https://doi.org/10.23974/ijol.2025.vol10.4.502}

Die University of Houston baute in den letzten Jahren ihre medizinischen
Ausbildungsgänge aus und richtete ein eigenes College of Medicine ein.
Der Text beschreibt nun, wie die Universitätsbibliothek ihren Bestand
auf die Interessen der neuen Studierenden und Forschenden ausrichtete.
Sie nutzte dazu eine Mischung aus der Auswertung verschiedener Daten und
weitergehender Überlegungen. Beispielsweise erstellte sie Listen der
Core-Zeitschriften der jeweiligen Unterfelder, die am College vertreten
sind, und überprüfte, zu welchen sie schon Zugang hatte oder zu welchen
sie Zugang schaffen müsste. Ein wenig ärgerlich am Text ist, dass sich
mehrfach von \enquote{traditionellen Methoden} der Bestandsbewertung
abgegrenzt wird, obgleich diese eigentlich auch hier nicht ersetzt,
sondern \enquote{nur} ergänzt werden. Aber offenbar ist es für die
diskursive Betonung der Innovationsfähigkeiten der Bibliothek wichtig.

Zudem handelt die Bibliothek selbstverständlich im US-amerikanischen
System von Hochschulbibliotheken, in denen die Kooperation lange nicht
so ausgeprägt ist wie in Europa. Beispielsweise werden
Konsortialverträge gar nicht erst angesprochen. Aber ansonsten liefert
der Text einmal einen Überblick darüber, wie in einer konkreten
Einrichtung eine Strategie für die Bestandsarbeit erstellt und dann auch
evaluiert wird. (ks)

\begin{center}\rule{0.5\linewidth}{0.5pt}\end{center}

Higgins, Colin (2026). \emph{The Continuing Relevance of DVDs and
Blu-ray Discs to Library Collections.} In: Library Resources \&
Technical Services 70 (2026) 1,
\url{https://doi.org/10.5860/lrts.70n1.8614}

Dieser Text argumentiert, dass Bibliotheken auch in Zukunft Filme auf
physischen Medien erwerben und im Bestand halten sollten. Es gibt, wie
hier dargestellt wird, einen klaren Trend dazu, Filme bei
Streamingdiensten zu lizenzieren. Aber -- so die wenig überraschende
Aussage im Artikel -- diese Dienste haben nicht den Fokus, Filme
langfristig zu erhalten, sondern andere ökonomische Perspektiven.
Gleichzeitig sind sie selber in ein System von befristeten
Lizenzierungen durch die Filmindustrie eingebunden, die es ihnen
unmöglich machen, solche langfristigen Perspektiven überhaupt
anzubieten. Nicht zuletzt gibt es eine ganze Reihe von Fällen, in denen
Filme bei diesen Streamingdiensten im Nachhinein verändert wurden. Das
passierte aus unterschiedlichen Gründen, ist aber selbstverständlich ein
Problem, wenn es darum geht, Filme als Kulturgut zu erhalten.

Der Text argumentiert, dass es die Aufgabe von Bibliotheken wäre, Filme
über die \enquote{Zufälligkeiten} dieses Streamingssystem hinaus zu
erhalten. Dies wäre nur möglich, wenn sie physisch vorhanden sind. Auch
das ist keine grundsätzlich überraschende Aussage. Interessant ist hier
aber, dass durch den Artikel Kriterien wie historische Verantwortung und
Diversität des Angebots für das Bestandsmanagement stark gemacht werden
und davor gewarnt wird, dabei zu kurzfristig -- also nur zu fragen:
\enquote{Was kann die Bibliothek sofort, für welchen Preis, mit welchem
Aufwand und zu welche Lizenzbedingungen ihren Nutzer*innen anbieten?} --
zu denken. (ks)

\subsection{2.4 Ukraine}\label{ukraine}

Der 2022 begonnene Krieg Russlands gegen die Ukraine hat auch zu einer
Reihe von Beiträgen in bibliothekarischen Zeitschriften geführt, von
denen hier einige in Englisch und Deutsch erschienenen referenziert
werden. Leider ist der Grund für dieses Interesse und auch für viele
Aktivitäten, die in diesen Beiträgen beschrieben werden, ein schlechter.
Es ist zu hoffen, dass dieser Krieg, wie alle Kriege, endlich beendet
wird. Gleichzeitig zeigen die Beiträge aber auch, dass Bibliotheken in
der Ukraine wichtige Arbeit leisten und dass sie dabei auf Solidarität
von Bibliotheken anderer Länder zurückgreifen können.

Dubrovina, Lyubov ; Kovtaniuk, Yurii ; Lobuzina, Kateryna ; Demianiuk,
Liudmyla (2024). \emph{The Vernadsky National Library of Ukraine in
Times of Independence and Martial Law: Development Strategy,
Preservation, and International Co-operation}. In: Bibliothek Forschung
und Praxis 48 (2024) 3: 639--647,
\url{https://doi.org/10.1515/bfp-2024-0062}

In diesem Artikel beschreiben Kolleg*innen aus der ukrainischen
Nationalbibliothek zuerst kurz die Geschichte der verschiedenen
Bibliotheken mit «nationalbibliothekarischem Auftrag» in ihrem Land.
Anschliessend stellen sie die Herausforderungen dar, welche durch den
Krieg für die Bibliotheken des Landes entstanden sind (Zerstörung,
Gefahren, aber auch Personalmangel), um danach auf die Strategien und
Arbeiten ihrer Bibliothek selber einzugehen.

Ziel der Arbeit der Bibliothek ist es aktuell, möglichst viele Medien,
Informationen und damit auch Kultur zu erhalten. Getan wird dies unter
anderem durch Digitalisierungsprojekte, bei denen die Digitalisate auch
in anderen Ländern gelagert werden, durch Sicherungsmassnahmen für die
Kulturgüter und durch die Integration in bestehende europäische
Infrastrukturen. Eine wichtige Rolle dabei spielt das Projekt «Saving
Ukrainian Cultural Heritage Online» (\url{https://www.sucho.org}). In
einem Kapitel wird auf die Unterstützung durch andere europäische
Nationalbibliotheken eingegangen, die in der Conference of European
National Libraries organisiert sind. Hervorgehoben werden dabei die
British Library, die Deutsche, Schwedische und Polnische
Nationalbibliothek. (ks)

\begin{center}\rule{0.5\linewidth}{0.5pt}\end{center}

Kulazhanka, Alexei (2025). \emph{ABD-Institutionen der Ukraine im Krieg:
Welche Schutzmassnahmen werden in einem Angriffskrieg angewendet, um
Bestände vor der Vernichtung zu retten?}. In: Informationswissenschaft:
Theorie, Methode und Praxis, 9 (2025): 255--268,
\url{https://doi.org/10.18755/iw.2025.15}

Der Artikel, wohl wie fast alle Beiträge dieser Zeitschrift entstanden
als Abschlussarbeit im gemeinsamen MAS-Studiengang
Informationswissenschaft der Universitäten Bern und Lausanne, fasst
zusammen, wie Archive, Bibliotheken und Dokumentationsstellen (die
ABD-Institutionen, wie sie im Titel genannt werden) in der Ukraine
vorgehen, um ihre Bestände zu schützen. Kulazhanka betont, dass «das
ordnungsgemässe Funktionieren von ABD-Institutionen {[}\ldots{]} eine
Normalität und Planbarkeit des Alltags {[}voraussetzt{]}», die im Krieg
nicht gegeben ist. Notfallpläne, wenn sie existieren, nähmen bewaffnete
Auseinandersetzungen immer nur als eine wenig wahrscheinliche
Katastrophe in den Blick, was dazu führt, dass ABD-Einrichtungen viel
weniger auf sie vorbereitet sind, als auf andere Krisensituationen. Dies
hätte auch für die Ukraine gegolten. Durch den Angriff Russlands hätten
sie neue Strategien entwickeln müssen, die im Artikel benannt werden
als: Evakuieren der Bestände, Sicherung vor Ort (beispielsweise in
Kellern) und Digitalisierung. Zu diesen drei Punkten werden jeweils
Beispiele aus der Ukraine genannt. (ks)

\begin{center}\rule{0.5\linewidth}{0.5pt}\end{center}

Gosart, Ulia ; Clair, Sara ; Swanson, Abbie (2025). \emph{The resistance
and resilience of Ukrainian public librarians}. In: IFLA Journal
{[}Online First{]}, \url{https://doi.org/10.1177/03400352251392112}
{[}Paywall{]}

Der kurze Text fasst Arbeiten zusammen, welche in ukrainischen
Öffentlichen Bibliotheken aufgrund des Krieges notwendig wurden. Dabei
werden diese Beispiele nicht diskutiert, sondern ein wenig wie
Nachrichten, die alle positiv oder «heldenhaft» klingen,
zusammengefasst. Zu erfahren ist, was alles getan wird, aber zum
Beispiel nicht, wie sich die Bibliothekar*innen dabei fühlen.

Zuerst werden einige Zahlen über den Zustand der Bibliotheken im Land
genannt: Im Mai 2025 waren 17\,\% von ihnen nicht mehr operationsfähig,
214 waren zerstört, fast 1.000 benötigten Reparaturen. Gleichzeitig
wurde wichtige Arbeit geleistet: Bibliotheken wurden zu Hilfszentren
ausgebaut, in denen Menschen Unterkunft und Hilfe finden konnten, oder
auch zu Luftschutzräumen. Es wurden Angebote im Bereich physische und
psychische Gesundheit aufgebaut, teilweise spezifisch für Kinder,
teilweise für die gesamte Bevölkerung. Es wurden aber auch Ausstellungen
organisiert (erwähnt werden solche in Zusammenarbeit mit der School of
Information, San Jose University in Kalifornien, woher die Autor*innen
stammen, was auch ein Beispiel für die Kooperation bibliothekarischer
Einrichtungen über Länder hinweg darstellt), bei denen Bilder gezeigt
wurden, mit denen ukrainische Kinder die Situation des Krieges
verarbeiteten. Zuletzt werden Beispiele genannt, in denen Bibliotheken
versuchen, Desinformation zu kontern, über die Desinformationsarbeit
Russlands aufzuklären und zudem den Krieg zu dokumentieren. (ks)

\begin{center}\rule{0.5\linewidth}{0.5pt}\end{center}

Johnston, Jamie ; Mierzecka, Anna ; Tóth, Máté ; Paul, Magdalena ;
Kisilowska-Szurmińska, Małgorzata ; Khosrowjerdi, Mahmood ; Vårheim,
Andreas ; Rydbeck, Kerstin ; Jochumsen, Henrik ; Hvenegaard Rasmussen,
Casper ; Ágústa, Pálsdóttir ; Olson, Anna ; Skare, Roswitha ;
Mathiasson, Mia Høj (2025). \emph{Public libraries' role in supporting
Ukrainian refugees: A focus on Hungary and Poland}. In: Journal of
Librarianship and Information Science, 57 (2024) 4:1141--1160,
\url{https://doi.org/10.1177/09610006241259490}

Vor allem basierend auf sechs Interviews (je drei in Polen und Ungarn,
dort je in einer gross-, mittel- und kleinstädtischen Bibliothek) gibt
dieser Text einen Überblick zur Arbeit von polnischen und ungarischen
Bibliotheken während der ersten Monate des Krieges Russlands gegen die
Ukraine. Beide Länder waren das Ziel von Personen, die aus der Ukraine
flüchteten; in beiden Ländern übernahmen neben der Zivilgesellschaft
auch Bibliotheken die Aufgabe, diesen Menschen Hilfe zu leisten. Wieder
in beiden Ländern musste diese Arbeit praktisch neu etabliert werden.
Während in anderen europäischen Ländern während der grossen
Migrationsbewegungen 2015 schon Erfahrungen dieser Art gesammelt wurden,
galt dies für Polen und Ungarn nicht. In diesen beiden Ländern gab es
2015 explizit eine Politik gegen die damals Flüchtenden, während die
2022 aus der Ukraine kommenden Personen sowohl von den Regierungen als
auch breiten Teilen der Gesellschaft unterstützt wurden.

Die Interviews zeigen zuerst, dass die Bibliotheken in beiden Ländern
sehr schnell auf die Fluchtbewegungen reagierten. In sehr kurzer Zeit
wurden sie zu Sammel- und Ausgabestellen für Spenden und zum
Aufenthaltsort für geflüchtete Personen. Dabei wurde betont, dass sich
hier professionelle und private Beweggründe mischten. Auf der einen
Seite sahen sich die Bibliotheken -- dieses Mal, aber offenbar kaum 2015
-- als «Ort für alle», auf der anderen Seite gab es eine Betroffenheit
der Bibliothekar*innen selber. In dieser Zeit standen Aufgaben im
Mittelpunkt, die -- so der Artikel -- sich nicht auf Medien bezogen.
Nach den ersten Monaten wandelte sich dann die Aufgaben der
Bibliotheken. Jetzt ging es eher darum, dass die Bibliotheken als
«Orte», die geflüchtete Personen aufsuchen konnten, als Einrichtung, an
denen Angebote für Kinder gemacht wurden und auch als Mediensammlung
wirkten. In dieser Zeit begannen auch die Bibliothekar*innen, die
geflüchteten Menschen in weitere Zielgruppen zu differenzieren und für
diese gesonderte Angebote zu unterbreiten, wobei sich vor allem ein
Fokus auf Kinder und Frauen ausprägte. Zudem wechselten die Aufgaben: Es
ging jetzt vor allem darum, Menschen das «Ankommen» im jeweiligen Land
zu ermöglichen, inklusive der Entwicklung von Sprachfähigkeiten und von
sozialen Verbindungen. Es gab, allerdings geringe, Unterschiede zwischen
der ungarischen und polnischen Situation, insbesondere, weil in Polen
mehr Menschen aus der Ukraine längerfristig verblieben als in Ungarn. In
Polen wurden die Bibliotheken so auch eher zu Einrichtungen, welche die
Integration förderten. (ks)

\section{3. Monographien und
Buchkapitel}\label{monographien-und-buchkapitel}

\subsection{3.1 Vermischte Themen}\label{vermischte-themen-1}

Campbell, Paul C. ; Nagle, Sarah (eds.) (2025). \emph{Librarians as
Researchers. Developing Our Scholarly Identities}. New York u.\,a.:
Bloomsbury. {[}gedruckt{]}

Wie so viele aktuelle Bücher aus dem Bibliothekswesen hat auch dieses
eine amerikanische \enquote{Herkunft}. Das macht direkte Vergleiche mit der Lage
in deutschen Bibliotheken etwas schwierig: Eine Unterscheidung von
Bibliotheksstellen nach \enquote{library positions} oder \enquote{faculty positions}
gibt es im deutschsprachigen Raum -- zumindest so weit ich weiß --
nicht, und zu den Anforderungen an \enquote{tenure track}-Stellen, wie zum
Beispiel publishing beziehungsweise research requirements, gibt es auch
kein direktes Pendant.

Nichtsdestotrotz kann dieser Sammelband dazu anregen, darüber
nachzudenken, was man \linebreak denn \enquote{erforschenswert} fände rund ums
Bibliothekswesen und wie man das angehen könnte.

In den Beiträgen gibt es anhand theoretischer Überlegungen und vieler
persönlicher Fallbeispiele Tipps zu grundlegenden Fragestellungen --
\enquote{What counts as scholarship for the purpose of one\textquotesingle s own
advancement versus scholarship that advances the discipline?} (S.\,10)
--, zur Erstellung einer \enquote{research agenda}, zum erfolgreichen Einstieg
ins Forschen zum Beispiel über Tagungs-Poster, zum Überwinden von Hürden
wie \enquote{time constraints and the assumptions others have about librarians}
(S.\,82) oder zu Formaten, die sich als gut praktikabel erwiesen haben,
wie Schreibgruppen, Communities of Practice und \enquote{Research Days}.

Online weiterlesen kann man danach dann zum Beispiel im ebenfalls in
einem Kapitel vorgestellten \enquote{Librarian Parlor}
(\href{https://deref-gmx.net/mail/client/8FHQ-s1Iumk/dereferrer/?redirectUrl=https\%3A\%2F\%2Flibparlor.com\%2F}{https://libparlor.com/})
und dem \enquote{LibParlor Online Learning} (LPOL)
(\href{https://deref-gmx.net/mail/client/oua9uIDXFEA/dereferrer/?redirectUrl=https\%3A\%2F\%2Flibparlor.com\%2Flpol-curriculum\%2F}{https://libparlor.com/lpol-curriculum/}).
(vv)

\begin{center}\rule{0.5\linewidth}{0.5pt}\end{center}

Syma, Carrye Kay ; Weiner, Robert G. ; Callender, Donell (Hrsg.) (2025).
\emph{Drawn to the stacks: essays on libraries, librarians and archives
in comics and graphic novels}. Jefferson, NC: McFarland.
\url{https://mcfarlandbooks.com/product/drawn-to-the-stacks/}
{[}gedruckt{]}

Wie werden Bibliotheken und Archive in Comics und Graphic Novels
dargestellt? Mit dieser Frage beschäftigen sich die 19 Beiträge in
diesem Band, nach eigenen Aussagen \enquote{the first academic volume to examine
the librarian and archival professional through the lens of sequential
art}.

Leider gibt es (vermutlich aus urheberrechtlichen Gründen) nur im
Artikel zur Reihe \enquote{xkcd} Abbildungen von Beispielen besprochener Comics
oder Cartoons. Das macht die Lektüre der Abhandlungen zu Comics, die man
nicht kennt, recht trocken. Dennoch können sie dazu einladen, sich die
besprochenen Werke mal anzusehen, um sich selbst ein Bild zu machen: Wer
sich für Comics und Bibliotheken interessiert, findet viele Anregungen
für Neuentdeckungen, und in bekannten Comics achtet man zukünftig
vielleicht noch mehr auf bibliothekarische oder archivarische Aspekte.

Sollte es mal einen zweiten Band dieser Art geben, wird der hoffentlich
noch weitere Werke berücksichtigen, die außerhalb der USA entstanden
sind, darunter zum Beispiel die bibliophilen Cartoons von Tom Gauld.
(vv)

\begin{center}\rule{0.5\linewidth}{0.5pt}\end{center}

Berger, Monica (2024). \emph{Predatory Publishing: and Global Scholarly
Communications}. (ACRL Publications in Librarianship, 81) Chicago:
Association of College and Research Libraries, 2024 {[}gedruckt{]}

Dieses Buch ist weit besser und fundierter, als man beim Thema Predatory
Publishing befürchten könnte. Es reproduziert gerade nicht einfach
Vorstellung von \enquote{betrügerischen Zeitschriften} und
\enquote{betrügerischen Forschenden} auf der einen, \enquote{ehrlichen
Verlagen} und \enquote{ehrlichen Forschenden} auf der anderen Seite,
sondern versucht, die komplexe Situation darzustellen. Zudem wird sich
mit den Problemen auseinandergesetzt, die vor gut zehn Jahren bei diesem
Thema existierten, wenn beispielsweise in \enquote{Beall's List}
Kriterien für \enquote{Predatory Publishing} angewendet wurden, die
Unterschiede zwischen Wissenschaftsstrukturen oder der
infrastrukturellen Ausstattungen von Hochschulen in verschiedenen
Ländern ignorierten und dazu führte, dass systematisch Zeitschriften aus
dem globalen Süden zu \enquote{predatory journals} erklärt wurden, weil
sie nicht Kriterien entsprachen, die in den USA, aber halt nicht überall
anders, sinnvoll waren. Zudem wird in diesem Buch hier explizit die
implizite Gleichsetzung von Open Access und Predatory Publishing
zurückgewiesen, welche gerade im Umfeld der genannten Liste vor gut zehn
Jahren vertreten wurde.

Vielmehr will das Buch systematisch das gesamte Themenfeld angehen. Es
bespricht nicht nur, warum Predatory Publishing für die Wissenschaft
schlecht ist, sondern auch, wie die verschiedenen Auffassungen davon,
was \enquote{ehrliches} oder \enquote{unehrliches} Publizieren ist
(sowohl zwischen verschiedenen Wissenschaftskulturen als auch innerhalb
solcher), entstehen und welche strukturellen Gründe vorliegen, die
Predatory Journals hervorbringen können. Gerade diese Gründe, hier unter
dem Schlagwort \enquote{Neoliberalismus} zusammengefasst, und nicht
unmoralisches Verhalten von einzelnen Forschenden wird als Motivation
für potentiell betrügerisches Verhalten angesehen. Zudem geht das Buch
explizit in einzelnen Kapiteln auf Diskussionen zu Predatory Publishing
und Situationen von Forschungseinrichtungen in verschiedenen Ländern
ein, verbleibt also gerade nicht beim US-amerikanischen Blick. Basis
dafür sind zahlreiche Forschungen zum Predatory Publishing, die in den
letzten Jahren durchgeführt wurden und hier sehr umfassend
zusammengeführt werden.

Im letzten Abschnitt des Buches wird versucht, praxisorientierte
Hilfestellungen für Bibliothekar*innen zu geben, die Forschende bei
Publikationsentscheidungen unterstützen wollen. Aber, da diese Situation
-- wie man im Buch selber lernen kann -- komplexer ist, als einfach das
Publikationswesen in \enquote{betrügerisch} oder \enquote{ehrlich}
einzuteilen, sind diese Hinweise auch recht umfangreich. Es geht immer
noch darum, Informationen über das Publikationswesen und seine
Auswirkungen weiterzugeben -- aber jetzt komplexere Informationen als
nur eine Liste von \enquote{schlechten Zeitschriften}, sondern eine
ganze Anzahl von bedenkenswerten Kriterien und Kontexten. (ks)

\begin{center}\rule{0.5\linewidth}{0.5pt}\end{center}

Mariátegui, José-Carlos (2023). \emph{Imagined Video Archives:
Strategies and Conditions of Video Art Collections in Latin America}.
In: Shtromberg, Elena ; Phillips, Glenn (edit.) (2023). \emph{Encounters
in Video Art in Latin America}. Los Angeles: Getty Research Institute,
2023: 91--109 {[}gedruckt{]}

Das Buch, zu welchem dieser Beitrag gehört, beschäftigt sich mit der
Geschichte der Videokunst in Lateinamerika. Das Buch selber geht von der
Prämisse aus, dass diese Kunst in Museen, Sammlungen und der
Kunstgeschichte massiv unterrepräsentiert ist, dass sie schlecht
erhalten wurde sowie dass sie sich zum grössten Teil gar nicht erhalten
hat. Dabei sei Videokunst gerade in den Militärdiktaturen während der
1970er und 1980er Jahre immer auch eine Form von Protest und Widerspruch
gegen diese Diktaturen gewesen und schon deshalb erinnerungswürdig.
Während sich andere Beiträge im Band mit einzelnen Festivals oder
Künstler*innen befassen, wurde im Beitrag von Mariátegui aufgearbeitet,
welche Archive und Sammlungen dieser Videokunst tatsächlich existieren.
Dabei kann kein vollständiger Überblick gegeben werden: So, wie die
damalige Videokunst oft lokal war, lokal produziert und gezeigt wurde,
und nur in einigen Fällen über diesen Rahmen hinaus wahrgenommen wurde,
ist es auch mit den Archiven. Viele sind Privatinitiativen, nur jeweils
einem kleinen Kreis von Eingeweihten bekannt.

Auf der Basis von Recherchen nach und Besuchen in solchen Archiven,
werden diese kategorisiert als (a) Sammlungen, die aus durchgeführten
Video Art Festivals entstanden, (b) Sammlungen in Institutionen wie
Museen, (c) persönlichen Sammlungen, insbesondere von
Videokünstler*innen selber und (d) Ausstellungsprojekten, für die
jeweils eine Sammlung zusammengetragen (und so zumindest dokumentiert)
wurden. Die Herausforderungen für alle diese Sammlungen bestehen im
Material selber, da die Kunst zumeist als Video oder Film vorliegt, an
technischen Herausforderungen, beispielsweise den verschiedenen Video-
und Filmformaten, und der Notwendigkeit, Abspielgeräte zur Verfügung
haben zu müssen sowie in vielen Fällen (aber nicht in allen Museen) in
der Finanzierung und Langlebigkeit der Initiativen selber. Der Überblick
zeigt zudem, dass trotz dieser Herausforderungen eine grosse Anzahl
solcher Sammlungen existiert, die oft von Engagierten mit grossem Elan
aufrecht erhalten werden. (ks)

\begin{center}\rule{0.5\linewidth}{0.5pt}\end{center}

Gilbert, Annette ; Hülhoff, Andreas (edit.) (2025). \emph{Library of
Artistic Print on Demand: Post-Digital Publishing in Times of Platform
Capitalism}. Leipzig: Spector Books, 2025 {[}gedruckt{]}

Dieses Buch basiert auf einer Sammlung von Artists Books, die nicht von
den Künstler*innen selber gedruckt, sondern über Print on
Demand-Services verbreitet wurden. Ausgangspunkt ist die Feststellung,
dass sich (auch) diese Form von Veröffentlichung in den letzten zehn
Jahren als künstlerische Praxis etabliert hat. In einem Projekt wurde
eine Sammlung von 100 dieser Bücher angelegt, welche anschliessend der
Bayerischen Staatsbibliothek in München gespendet wurde (die schon eine
eigene Sammlung von Artists Books besitzt). Teil dieser Publikation ist
ein bebilderter und kommentierter Katalog dieser Sammlung.

Rund zwei Drittel dieses schweren Buches (A4, knapp über 700 Seiten)
besteht aber aus Essays, welche diese Publikationspraxis dokumentieren,
diskutieren und kritisieren. Grundsätzlich \linebreak wird dabei immer wieder der
Widerspruch erwähnt, dass hier Künstler*innen oft mit einem
gesellschaftlich kritischen Blick Bücher über Plattformen vertreiben,
die vor allem von grossen Firmen, insbesondere Amazon, betrieben werden.
Einige thematisieren dies in ihren eigenen Büchern, einige wägen die
Vor- und Nachteile ab, einige besprechen auch, dass alternative Formen
von Print on Demand in der Vergangenheit gescheitert sind. Die
Diskussion über den Widerspruch kommt nie an ein Ende.

Zusammengefasst bietet das Buch einen Meilenstein dieser künstlerischen
Praxis, da hier eine offenbar lebendige \enquote{Szene} das erste Mal
zusammengeholt und dargestellt wird (selbstverständlich ohne Anspruch
auf Vollständigkeit, wie schon im Vorwort betont wird, schon weil gerade
Bücher aus Ländern des Globalen Südens fehlen). (ks)

\subsection{3.2 Bibliotheksmanagement}\label{bibliotheksmanagement}

Hunt, Amber ; Ruane, Elizabeth ; Sopka, Stephanie (edit.) (2024).
\emph{Closing a College Library}. Chicago: Association of College and
Research Libraries, 2024 {[}gedruckt{]}

Das US-amerikanische Hochschulsystem ist geprägt von der Existenz vieler
kleiner, privater (oft auch von religiösen Einrichtungen getragenen)
Hochschulen. Viele dieser kleinen Einrichtungen, insbesondere die den
Fachhochschulen vergleichbaren Colleges, stehen unter starkem
finanziellen Druck, da die Studierendenzahlen massiv zurückgehen.
Offenbar immer öfter werden sie geschlossen, was heissen kann, dass sie
mit grösseren Einrichtungen fusionieren oder aber, dass sie tatsächlich
vollständig aufgelöst werden. Im Rahmen dieser Schliessungen werden auch
die Bibliotheken dieser Einrichtungen geschlossen. In diesem Buch
versammelt sind Erfahrungsberichte von einigen Bibliotheken, mit denen
dies in den letzten Jahren -- zufällig immer in der ersten Hochzeit der
COVID-19-Pandemie -- tatsächlich passierte. Allerdings ist auch
auffällig, dass hier sehr die US-amerikanischen Strukturen eine Rolle
spielen. So handelt es sich immer um private Colleges, die dann, wenn
sie fusionieren, auch immer mit anderen privaten Hochschulen
zusammengeschlossen wurden. Immer geht es um die Verwaltung von Spenden
(Endowments) und Special Collections sowie Hochschularchive. Zudem geht
es immer wieder um Geld. Beispielsweise wurde bei einer Bibliothek der
Transfer von Büchern schwierig, weil diese inklusive der Gebäude und des
Geländes verkauft werden sollten. Zudem werden auch immer wieder
Herausforderungen durch verschiedene Konsortien angesprochen (da es,
anders als in Europa heute die Norm, kein nationales Konsortium gibt,
sondern ständig wechselnde Konsortialstrukturen).

Oder anders: Es ist schwer, dass alles direkt auf die Bibliotheken im
DACH-Raum zu übertragen. Zwar stehen auch hier die kleinen Hochschulen
vor dem Problem sinkender Studierendenzahlen, aber dabei handelt es sich
fast immer um staatliche Einrichtungen, deren Bibliotheken auch mehr
vernetzt sind als in den USA. Bisher wurden auch noch sehr wenige von
ihnen geschlossen. Wenn, dann geht es im DACH-Raum wohl vor allem um das
Auflösen von Teilbereichsbibliotheken eines
Hochschulbibliothkekssystems, was eine sehr andere Problemlage ist.

Was aus den Beispielen aber zu lernen ist, ist, (a) dass so bald möglich
nach dem Bekanntwerden der Schliessung eine Planung stattfinden sollte,
die wenn möglich auch Bibliotheken und Archive einbeziehen sollte,
welche Bestände übernehmen können, (b) dass eine schnelle
Inventarisierung sowie Sicherung von möglichst vielen Daten und
Metadaten sinnvoll ist, (c) dass sich oft Personen eigenhändig an den
Beständen der zu schliessenden Bibliotheken \enquote{bedienen}, wenn dem
nicht schnell ein Riegel vorgeschoben wird. Zudem ist dem Schlusswort zu
entnehmen, dass die Schliessungswelle von Colleges und Bibliotheken in
den USA noch lange nicht an ein Ende gelangt ist. (ks)

\begin{center}\rule{0.5\linewidth}{0.5pt}\end{center}

Hirsch, Sandra (edit.) (2024). \emph{Library 2035: Imagining the Next
Generation of Libraries}. Lanham ; Boulder ; New York ; London: Rowman
\& Littlefield, 2024 {[}gedruckt{]}

Das Interessanteste an diesem Buch ist, dass es aus der Zeit gefallen
scheint. Vor zehn bis fünfzehn Jahren erschienen ständig Bücher und
Artikel, in denen vor allem Bibliothekar*innen oder Forschende aus der
Bibliothekswissenschaft Aussagen darüber machten, wie die Bibliothek der
Zukunft aussehen wird. Eine ganze Anzahl von Studien und Trendberichten
versuchte damals auch, diese Frage systematischer anzugehen. Oft wurden
diese Anstrengungen mit stattfindenden oder vorhersehbaren Trends
begründet, die auch auf Bibliotheken wirken würden. Teilweise wurden
damals weitgreifende Vorhersagen gemacht und gewarnt, dass Bibliotheken
sich aktiv wandeln würden müssen, wenn sie nicht rapide an Bedeutung
verlieren wollten. Aber gleichzeitig wurde ihnen damals auch schon eine
bleibende Bedeutung zugeschrieben: Für das Management digitaler Medien
und Daten sowie für die Pflege von Communities.

Heute ist weniger Zukunft, könnte man sagen. Immer weniger Kolleg*innen
versuchen sich an solchen Vorhersagen. (Und die, die es tun -- das eine
subjektive Wertung des Rezensenten --, klingen manchmal so, als würden
sie einfach seit zehn oder fünfzehn Jahren immer wieder das Gleiche
sagen, was auch nicht so richtig nach \enquote{Zukunft} klingt.) Forecasts sind
heute weniger beliebt und wenn, dann eher von negativer Sicht geprägt.

Dieses Buch jetzt scheint noch einmal den Elan, der sich 2010--2015 im
(englisch- und deutschsprachigen) Bibliothekswesen finden liess,
aufzurufen. Aber es liest sich auch ein wenig wie ein zu spät kommender
Beitrag, der nicht mehr so richtig in die Landschaft passt. Auch wenn
nicht so richtig klar ist, warum eigentlich nicht: Als wäre es falsch,
positive Zukunftsentwürfe zu liefern -- was eigentlich nie der Fall ist,
auch wenn der Zeitgeist negativer gestimmt scheint.

Die Beiträge hier, alle aus dem US-amerikanischen Raum, sind kurz, nie
länger als zehn Seiten. In ihnen nennen Bibliothekar*innen und
Bibliothekswissenschaftler*innen je kurz ein Thema wie Nachhaltigkeit,
Diversity, Virtual Technologies oder die Entwicklung des Arbeitsmarktes
und machen dann kurze, nie so richtig fundierte, Aussagen darüber, wie
sich diese im Bezug auf Bibliotheken (wieder unausgesprochen aber
offensichtlich denen in den USA) verändern werden. Am Ende betonen sie
fast in jedem Fazit, dass Bibliotheken sich entscheiden müssten, welche
Zukunft sie anstreben und wie sie dorthin kommen wollen. Das ist alles
auf der einen Seite ganz aufbauend -- die Autor*innen gehen offenbar
immer davon aus, dass Bibliotheken die aktuellen Herausforderungen
meistern werden --, aber auch irgendwie wenig glaubwürdig. Die Beiträge
lesen sich weniger wie überzeugende (oder gar in Daten sichtbare)
Zukünfte der Bibliotheken und viel mehr wie Texte, in denen sich
Bibliothekar*innen gegenseitig Mut machen. (ks)

\begin{center}\rule{0.5\linewidth}{0.5pt}\end{center}

Smith, Christine F. (2025). \emph{Platform Power and Libraries}. (Series
on Critical Information Organisation in LIS, 3) Sacramento: Library
Juice Press, 2025 {[}gedruckt{]}

Hier werden vier kurze Beiträge und eine längere Einleitung versammelt,
welche den Stand von \enquote{Platformisation} im Bibliotheksbereich
beschreiben wollen. (Die Beiträge stammen aus den USA, Kanada sowie
Australien und stehen im jeweiligen Kontext.) Mit Platformisation ist
das Entstehen von Intermediären gemeint, die zwischen Medienanbietern
(Verlagen, Filmdistributionen, Labels et cetera) und Bibliotheken
stehen, die sich wiederum potentiell immer mehr an Plattformen wie
Spotify, Netflix oder den zahlreichen Video-Streaming-Diensten
orientieren. Obgleich die meisten Anbieter für Bibliotheken
grundsätzlich ein Interesse daran hätten, Bibliotheken zu bedienen,
würden sie immer mehr auch von der Idee getrieben, Profit mit der
Auswertung von Daten und der Skalierung ihrer Angebote zu machen.

Die Einleitung stellt all dies in den Kontext einer Kritik
kapitalistischer Strukturen. Auch in den Beiträgen wird mehrfach
erwähnt, dass das Handeln der Intermediäre durch diese Strukturen zu
erklären ist (also dadurch, dass Firmen einen Profit anstreben müssen
und dass sie ständig an der Ausweitung ihres Geschäftes interessiert
sind). Bibliotheken sollten dies mit bedenken und nicht in
Verschwörungstheorien verfallen oder den Intermediären unterstellen,
diese wollten aktiv Bibliotheken schädigen. Gleichwohl beschreiben die
Texte die Situationen als schwierig: Bibliotheken würden bei digitalen
Medien immer mehr die Möglichkeit verlieren, den eigenen Bestand
auszuwählen und zu managen. Ausserdem wäre es immer schwieriger, für die
langfristige Archivierung zu sorgen. Der Grossteil der Texte beschreibt
diese Situation und betont, dass Bibliotheken darüber miteinander reden
müssten, bietet aber keine Lösungsansätze. Ein wenig überraschend ist
dies, weil sich eigentlich aufdrängt, darüber nachzudenken, ob (a)
Bibliotheken ihre Verhandlungsmacht gegenüber den Intermediären durch
Konsortien stärken oder (b) selber genossenschaftlich Intermediäre
schaffen und unterstützen sollten. Ansonsten ist beim Lesen der Beiträge
nicht so richtig klar, warum die kapitalismuskritische Perspektive in
der Einleitung überhaupt so stark betont wurde. In den Beiträgen wird
sich nicht auf sie bezogen. (ks)

\subsection{3.3 Diversity und
Bestandsmanagement}\label{diversity-und-bestandsmanagement}

Barber, Erika ; Bauder, Julai ; Behounek, Micki ; Jones, Chris ; Reed,
Kayla ; Rodrigues, Elizabeth (edit.) (2024). \emph{Supporting Diversity
through Collection Evaluation, Development, and Weeding}. (CLIPP, 48).
Chicago: Association of College and Research Libraries, 2024
{[}gedruckt{]}

Dieses Buch erschien in einer Reihe, welche praxisorientierte und recht
kurze Publikationen für die US-amerikanischen Hochschulbibliotheken
veröffentlicht. (CLIPP steht dabei für \enquote{College Library
Information on Policy and Practice}.) Es verspricht im Titel einen
Überblick zu geben zu Möglichkeiten und Tools, um das Bestandsmanagement
daraufhin auszurichten, dass der Bestand selber diverser wird. Im
einleitenden Text wird dann darauf eingegangen, dass
\enquote{Diversität} ein Begriff ist, der für verschiedene Dimensionen
gelten kann: für den Inhalt von Medien, für die Identitäten und
Herkünfte der Autor*innen et cetera. (Nicht erwähnt wird
überraschenderweise eine mögliche Diversität von Verlagen und anderen
herausgebenden Körperschaften.)

Selbstverständlich ist das Thema relevant und es wäre zu begrüssen, wenn
dieses Buch tatsächlich eine Handreichung zur Verfügung stellen würde.
Was es allerdings hauptsächlich ist, ist die Darstellung der Ergebnisse
einer Umfrage zu diesem Thema unter US-amerikanischen
Hochschulbibliotheken -- aus der man erfährt, dass das Thema trotz aller
politischen Versuche, die dagegen stehen, Relevanz hat und auch in
vielen Einrichtungen Teil der Bestandsarbeit ist, aber das gleichzeitig
nicht alle Bibliotheken grundsätzlich so professionell arbeiten, wie man
vielleicht erwarten würde -- und dann eine ganze Sammlung von Collection
Development Policies von Bibliotheken, welche Diversity-Ziele für ihren
Bestand gesetzt haben oder in ihnen auch festgelegt haben, welche Tools
sie einsetzen, um diese zu erreichen. Aber das ist dann immer nur ein
Teil dieser Dokumente, die jeweils in ihrer Gesamtheit abgedruckt
werden. Von den rund 120 Seiten nehmen sie fast 80 ein. Anders gesagt:
Es ist eigentlich ein kurzer Artikel mit einem langen Anhang, der zu
einem Buch \enquote{erhoben} wurde.

Eine interessante Erkenntnis für Bibliotheken aus dem DACH-Raum ist
vielleicht noch, dass eine ganze Anzahl von US-amerikanischen
Bibliotheksausstattern Tools anbietet, damit dortige Bibliotheken ihren
Bestand auf Diversity-Ziele hin analysieren können. Aber leider werden
diese nur kurz in einem Satz erwähnt und nicht weiter dargestellt. (ks)

\begin{center}\rule{0.5\linewidth}{0.5pt}\end{center}

Kohn, Karen (2025). \emph{Assessing Academic Library Collections for
Diversity, Equity, and Inclusion}. New York: Bloomsbury Inc., 2025
{[}gedruckt{]}

Den Anspruch, welchen das zuvor besprochene Buch erhebt -- nämlich
Möglichkeiten für die Umsetzung von Diversity-Zielen im
Bestandsmanagement von Hochschulbibliotheken aufzuzeigen -- erfüllt
dieses Buch viel besser. Es führt sehr direkt in das Thema ein, geht
dann, zuerst zusammengefasst, dann noch ausführlicher in je einem
Kapitel, auf vier Ansätze (List-Checking, Metadata Searching, Diversity
Coding und Assessing Institutional Efforts) ein, welche direkt im
Management des Bestands umgesetzt werden können. Zu drei der Ansätze
werden zudem in je einem weiteren Kapitel konkrete Use cases aus
US-amerikanischen Bibliotheken gezeigt (nur die Arbeit mit Metadaten
wird nicht gesondert vorgeführt). Die Vor- und Nachteile der Ansätze
werden besprochen und es wird auch gezeigt, welche für kurzfristige
Projekte und welche vor allem für langfristige Wirkungen geeignet sind.
Das alles ist in einem eher unaufgeregten Ton gehalten und zudem so
dargestellt, dass es in andere Bibliotheken übertragbar wird.

Sichtbar wird allerdings auch, dass es in den USA mehr Anbieter für
Medien gibt, welche diese Arbeit der Bibliotheken schon unterstützen
(indem sie beispielsweise explizit diverse Medien als Paket anbieten
oder aber, indem sie sich auf bestimmte Themen wie afroamerikanische
Geschichte spezialisieren) sowie Projekte wie kontinuierlich gepflegte
Bibliographien von Literatur von Autor*innen aus marginalisierten
Gruppen, auf die bei der Bestandsauswahl zurückgegriffen werden kann.
Kohn stellt dar, wie deren Angebote von Bibliotheken genutzt werden
können. Aber im DACH-Raum ist dies, bislang zumindest, wegen fehlender
Anbieter oder Projekte kaum reproduzierbar. Trotzdem ist dieses Buch die
bessere Wahl, wenn sich über praxisorientierte Möglichkeiten des
diversitygerechten Bestandsmanagements informiert werden soll. (ks)

\subsection{3.4 Bibliotheks- und
Mediengeschichte}\label{bibliotheks--und-mediengeschichte}

Schmidt, Siegfried (2024). \emph{Das} Hirtenwort der deutschen Bischöfe
zur katholischen Literatur der Gegenwart \emph{aus dem Jahre 1955:
Inhalt, Entstehung und Wirkungsgeschichte.} In: Schmidt, Siegfried ;
Reudenbach, Hermann-Josef (Hrsg.). Katholische Lesekultur im Rheinland:
von den \emph{Rheinischen Volksblättern} zum Hirtenbrief über
Gegenwartsliteratur 1955. Vorträge auf dem Vierten Kolloquium zum
Katholischen Buch- und Verlagswesen im Rheinland am 17. November 2023.
(Libelli Rhenani, 86). Köln: Erzbischöflische Diözesan- und
Dombibliothek mit Bibliothek St. Albertus Magnus, 2024, S. 15--104
{[}gedruckt{]}

Im Jahre 1955 veröffentlichten (vorgeblich) alle deutschen katholischen
Bischöfe ein Hirtenwort dazu, wie die katholische Öffentlichkeit sich in
Bezug auf die damalige Literatur verhalten sollte -- und damit waren
auch die Nutzer*innen der damals noch direkt an diese Öffentlichkeit
gerichteten katholischen Büchereien, die im Borromäusverein organisiert
waren, sowie die dort tätigen Bibliothekar*innen gemeint. Ein Hirtenwort
wird normalerweise von den Pfarrern in den Kirchen von der Kanzel
verlesen. Bei diesem gab es aber, wie der Artikel von Schmidt zeigt,
verschiedene Handhabungen. Meist wurde es nur in den katholischen
Amtsblättern veröffentlicht.

Das Hirtenwort (im Folgenden wird hier aus der Version im Anhang des
Artikels von Schmidt (89--94) zitiert) betonte explizit, dass es Themen
gab, die in der Gegenwartsliteratur besprochen wurden, die zu lesen aber
für die katholischen Laien nicht zu empfehlen sei -- oder aber explizit
verboten waren. (Letzteres war ernst gemeint: Der Index librorum
prohibitorum der katholischen Kirche galt noch bis zum Zweiten
Vatikanischen Konzil, das 1965 endete.) Genannt wurden von diesen Themen
explizit \enquote{anormale Veranlagungen und Verhaltensweisen} im
\enquote{Geschlechtlichen des Menschen} (was offensichtlich alles von
ausserehelichem Sex und bis zu allen Formen der sexuellen Identität
ausserhalb der heteronormativen, cis-gendered Norm meinte, wie man es
heute ausdrücken würde), Selbstmord, \enquote{sittliche Dekadenz} und
die Gefahr, die Lösung von persönlichen Problemen nicht in
\enquote{Gottes barmherzige{[}r{]} Gnade} zu suchen. Im Hirtenwort
wurden die Pfarrer ermahnt, auf die Lektüre der katholischen Laien zu
achten und dann letztere direkt angesprochen: \enquote{Unsere Gläubigen
aber seien ermahnt, bei der Wahl ihrer Lektüre zu prüfen, was ihnen
zuträglich ist, was ihnen echte Anregung und wirksame Förderung
vermittelt, was ihr Christentum lebendiger, reifer und reiner macht. Vor
allem möge man bedenken, daß nicht für jeden jedes Buch frommt und
namentlich die Kinder und die Jugend vor bedenklichen Einflüssen und
verfrühten Begegnungen mit problematischer Literatur bewahrt werden
müssen.}

Das Hirtenwort selber, welches sich heute ein wenig wie eine Satire
darüber liest, wie man sich das Denken wertender, konservativer
Kirchenkreise vorstellt, war selbstverständlich in der Forschung
bekannt. Schmidt zeigt in seinem umfangreichen, mit mehreren Anhängen
ausgestatteten Beitrag jetzt aber, unter anderem auf Basis von
Archivalien, die Entstehung und Wirkung des Textes auf. Ein Hauptpunkt
ist dabei, dass der Text vor allem von einem Erzbischof (Josef Kardinal
Frings, Köln) und einem von ihm geförderten Priester (Dr.~Leo Koep), der
kurz vorher als Generalsekretär beim Borromäusverein tätig geworden war
(1966 übernahm er dann die Leitung des Vereins), erstellt und seine
Verbreitung vorangetrieben wurde. Die anderen Bischöfe wurden
überrumpelt, indem Frings praktisch eine Deadline setzte, damit das
Hirtenwort noch vor Weihnachten verbreitet werden könne, da es sonst
angeblich keine Wirkung hätte. Das erklärt auch, warum das Hirtenwort in
den verschiedenen Kirchenprovinzen unterschiedlich gehandhabt wurde:
Teilweise wurde es nur publiziert, teilweise mit einordnenden Vorworten
versehen, teilweise wurde es Pfarrern empfohlen, es in der Kirche zu
verlesen, teilweise auch nicht. In Briefen an Frings äusserten viele
Bischöfe grundsätzlich, dass die Hauptaussagen richtig wären, aber man
über den konkreten Text noch hätte diskutieren müssen.

Als Antrieb zum Entstehen des Hirtenworts stellt Schmidt zwei
Entwicklungen heraus: Auf der einen Seite gab es in der katholischen
Literatur der damaligen Zeit einen Aufschwung, bei dem sich in
literarischer Form den realen Problemen der Menschen angenommen wurde
(die sich in einigen dieser Werke auch nicht auf die \enquote{richige
katholische Weise} lösten, sondern beispielsweise in einem Selbstmord
enden konnten), wenn auch immer aus konservativer katholischer Sicht
geschrieben. (Bekannt ist diese Richtung unter den Bezeichnungen
\enquote{renouveau catholique} oder \enquote{Katholische Erneuerung}).
Frings sah diese Literatur als Gefahr an, oder genauer, dass
\enquote{einfache Laien} diese lesen würden. Auf der anderen Seite
hatten die französischen Bischöfe Anfang 1955 eine ähnliche
Verlautbarung veröffentlicht (\enquote{Directives de l'Assemblée des
Cardinaux et Archevêques au sujet des lectures}), die explizit eine
Lenkung der Literatur der katholischen Laien propagierte (mit dem
erklärten Ziel, ihnen zu helfen, \enquote{katholische Lösungen} ihrer
eigenen Probleme zu finden). Schmidt wertet diese als Anstoss, aber
nicht Hauptgrund für das Hirtenwort.

Im Ganzen zeigt sich, dass die im Hirtenwort vertretene Position damals
tatsächlich in der katholischen Kirche vertreten wurde, aber dass es
auch unter den Bischöfen schon damals teilweise ein grosses Unbehagen an
ihr gab. (Nach dem Zweiten Vatikanischen Konzil oder gar heute ist diese
Position gänzlich unhaltbar. Kirche und Borromäusverein unterstützen
vielmehr das Recht auf die freie Lektürewahl und -interpretation.) Der
Beitrag von Schmidt zeigt aber, dass diese Position auch nicht aus einer
allzu fernen Vergangenheit stammt. (ks)

\begin{center}\rule{0.5\linewidth}{0.5pt}\end{center}

Schmidt, Siegfried ; Reudenbach, Hermann-Josef (Hrsg.) (2018).
\emph{Beiträge zum Katholischen Buch- und Verlagswesen im Rheinland im
19. und 20. Jahrhundert}. (Libelli Rhenani, 72) Köln: Erzbischöfliche
Diözesan- und Dombibliothek mit Bibliothek St.~Albertus Magnus, 2018
{[}gedruckt{]}

Schmidt, Siegfried ; Reudenbach, Hermann-Josef (Hrsg.) (2020).
\emph{Neue Beiträge zum Katholischen Buch- und Verlagswesen im Rheinland
im 19. und 20. Jahrhundert}. (Libelli Rhenani, 78) Köln: Erzbischöfliche
Diözesan- und Dombibliothek mit Bibliothek St.~Albertus Magnus, 2020
{[}gedruckt{]}

Schmidt, Siegfried ; Reudenbach, Hermann-Josef (Hrsg.) (2022).
\emph{Kleinschriften} -- \emph{Periodica} -- \emph{Prachtwerk} --
\emph{Predigtliteratur: Vorträge auf dem Dritten Kolloquium zum
Katholischen Buch- und Verlagswesen im Rheinland am 18. März 2022}.
(Libelli Rhenani, 83) Köln: Erzbischöfliche Diözesan- und Dombibliothek
mit Bibliothek St.~Albertus Magnus, 2022 {[}gedruckt{]}

Schmidt, Siegfried ; Reudenbach, Hermann-Josef (Hrsg.) (2024).
\emph{Katholische Lesekultur im Rheinland: von den} Rheinischen
Volksblättern \emph{zum Hirtenbrief über Gegenwartsliteratur 1955.
Vorträge auf dem Vierten Kolloquium zum Katholischen Buch- und
Verlagswesen im Rheinland am 17. November 2023}. (Libelli Rhenani, 86).
Köln: Erzbischöfliche Diözesan- und Dombibliothek mit Bibliothek
St.~Albertus Magnus, 2024 {[}gedruckt{]}

Seit 2017 sind, in nicht ganz regelmässigen Abständen, die unter anderem
durch die Covid-19-Pandemie bedingt waren, in Köln vier Kolloquien zum
Thema Katholisches Buch- und Verlagswesen im Rheinland im 19. und 20.
Jahrhundert durchgeführt worden. Anwesend waren immer rund 20 Personen,
organisiert wurden sie in der Erzbischöflichen Diözesan- und
Dombibliothek und zwar offenbar federführend von jetzt zwei ehemaligen
Leitern dieser Bibliothek (einer ist im Ruhestand, der andere offenbar
als Priester anderweitig verpflichtet), die auf den Kolloquien auch oft
vortrugen und als Herausgeber der Bände fungieren. Diese vier Bücher
stellen die Tagungsbände dar, die auch je die Teilnehmenden und die
Tagungsprogramme dokumentieren und Platz für den Abdruck Vorträgen
bieten, die wegen Krankheit nicht gehalten wurden. (Ein Beitrag aus dem
zuletzt erschienenen Band wurde weiter oben gesondert vorgestellt.) Für
die Bibliotheksgeschichte relevant ist die Kolloquiumsreihe dadurch,
dass in Bonn seit 1845 der Borromäusverein sitzt, welcher (bis heute)
als Aufgabe vor allem die Förderung und Anleitung katholischer
Büchereien hat. Offenbar existieren bis heute zahlreiche Archivalien zu
diesem Verein, welche im Laufe der Kolloquien herangezogen wurden.

Das 19. und 20. Jahrhundert stellte für das katholische Milieu die Zeit
dar, in denen es sich praktisch erst als eigenständiges, geschlossenes
Milieu \enquote{fand}, welches nicht einfach durch spezifische religiöse
Glaubensgrundsätze und -praktiken gekennzeichnet war, sondern praktisch
einen eigenen Lebensraum bildete, in denen die Mitglieder (das waren
alle Katholik*innen, nicht einfach die kirchlichen Institutionen und
deren Mitglieder) praktisch ihr gesamtes Leben führen können sollten.
Für die Medien des Milieus hiess dies, dass praktisch alle Formen der
jeweils aktuellen Medien aufgegriffen wurden (beispielsweise
Kleinschriften, Zeitschriften oder Tageszeitungen), aber explizit für
Katholik*innen gedacht sowie mit eigener Ästhetik gestaltet, und das
gleichzeitig eigene, spezifisch katholische Medienformen (beispielsweise
Wallfahrtsbücher oder Predigtensammlungen) etabliert wurden. Gerade in
letzterem war das Milieu unbestreitbar innovativ. Mit der Zeit formierte
sich erst das Milieu, spielte dann wichtige Rollen in der Geschichte (im
Kolloquium explizit thematisiert ist die rheinische und die deutsche
Geschichte) und löste sich dann, so wie andere Milieus, ab dem 1960er
Jahren als eigenständiges, geschlossenes Milieu auf. Es war,
medienhistorisch gesprochen, eine bewegte Zeit.

Die Beiträge in den vier Bänden beschäftigen sich nun mit sehr
unterschiedlichen Teilaspekten dieser Geschichte. Beispielsweise werden
einzelne katholische Verlage oder Autor*innen untersucht. Oder es wird
auf die Gestaltung von Büchern bestimmter Verlage eingegangen. Mehrfach
werden Zeitschriften und ihr Kontext vorgestellt (beispielsweise eigene
Zeitschriften für Ministranten im 20. Jahrhundert oder ein
Literaturblatt, welches zum Sammelpunkt einer eigenen theologischen
Strömung wurde, die sich dann als Altkatholische Kirche abspaltete).
Auch dem Borromäusverein werden mehrfach Beiträge gewidmet (alle von
Siegfried Schmidt), zudem taucht er in anderen Beiträgen auch an
unerwarteter Stelle immer wieder auf. Sichtbar wird, dass er als Verein
eine grosse Reichweite hatte, die sich auch dadurch erklärt, dass die
katholischen Büchereien -- zumeist ehrenamtlich betrieben und in hoher
Zahl vorhanden -- einen direkten Einfluss auf das Leben einer grossen
Zahl von katholischen Laien hatten. Allerdings steht bei diesen
Beiträgen die eigentliche bibliothekarische Arbeit -- also, was in den
Büchereien tatsächlich getan wurde, wie sie organisiert waren, wie die
Medien ausgewählt wurden und so weiter -- nie im Fokus, sondern immer
wieder andere Themen (zum Beispiel Kleinschriften, die der Verein für
wenige Jahren nach 1945 herausgab, oder eine Buchgemeinschaft, die er
betrieb).

Leicht kommt der Eindruck auf, dass die Kolloquien für die Autor*innen
der Beiträge eine Möglichkeit darstellten, ein sie interessierendes,
spezielles Thema teilweise sehr ausführlich und intensiv darzustellen.
Beiträge von über 50, manchmal über 100 Seiten, mit längeren Anhängen,
sind keine Seltenheit. Das heisst nicht, dass sie uninteressant wären,
aber teilweise gleichen sie ausführlichen wissenschaftlichen Beiträgen
(wobei der wissenschaftliche Hintergrund bei Beitragenden auch immer
wieder sichtbar wird, wenn beispielsweise die Fussnoten den grössten
Teil einer Druckseite einnehmen), und gerade nicht Vorträgen.
Gleichzeitig vereinen die Beiträge die Vor- und Nachteile einer
spezifischen Regionalgeschichtsschreibung. Auf der einen Seite schätzen
die Autor*innen ihre Themen und die beschriebenen Personen wert.
Vermutlich (bei einigen, beispielsweise Priestern auch explizit) sind
sie selber Katholik*innen, welche es deshalb vielleicht weniger
erstaunlich finden als andere Menschen, wenn sich Autor*innen des frühen
20. Jahrhunderts Gedanken um das Wunderwirken von Heiligen machten oder
wenn in Broschüren diskutiert wurde, wie Mütter ihren Kindern das
richtige Beten beibringen könnten. (Der Rezent als Atheist im 21.
Jahrhundert muss zugeben, dass ihm eine solche Wertschätzung nicht bei
allen Themen möglich gewesen wäre.) Das ist hilfreich, da so besser
beschrieben wird, was im katholischen Milieu damals gedacht und gelebt
wurde. Sie äussern dabei auch Kritik, sind also nicht apologetisch, aber
akzeptieren auch bestimmte Denkweisen als zeittypisch, die heute wenig
nachvollziehbar erscheinen. Auf der anderen Seite verfangen sich die
Beiträge teilweise sehr in spezifischen Kleinfragen oder verweisen auf
Positionen, Entwicklungen oder lokale Gegebenheiten, die für
Katholik*innen aus dem Rheinland vielleicht bekannt sind, welche aber
für andere Leser*innen immer wieder Nachrecherchen nötig machen.

Alles in allem verströmen die Bände den Habitus des Zusammentreffens von
Personen, die an einer spezifischen Mediengeschichte interessiert sind
und die hier einen Ort gefunden haben, um ihre Interessen auszuleben.
Das ist zu begrüssen, zumal das auf den Treffen versammelte Wissen durch
die Bände ja auch zugänglich gemacht wird. Zu wünschen ist -- ausser,
dass diese Reihe in Zukunft weitergehen wird -- wenn überhaupt, dass
erstens in Zukunft öfter Bezug genommen wird auf die Mediengeschichte
ausserhalb des katholischen Milieus zur damaligen Zeit. Es wird nämlich
an vielen Stellen nicht klar, ob die beschriebenen Entwicklungen
spezifisch waren oder Teil der gesamten Entwicklung von Medien und
Gesellschaft. Und zweitens wären Beiträge, die auf die bibliothekarische
Arbeit des Borromäusvereins eingehen, ein Desiderat, das geschlossen
werden sollte. (ks)

\begin{center}\rule{0.5\linewidth}{0.5pt}\end{center}

Haver, Gianni (2018). \emph{La presse illustrée: Une histoire romande}.
(Savoir Suisse) Lausanne: Presses polytechniques et universitaires
romandes, 2018 {[}gedruckt{]}

Savoir Suisse ist nominell ein eigener Verlag, praktisch aber als
Imprint unterhalb des Universitätsverlages der schweizerischen
Bundeshochschule EPFL in Lausanne angesiedelt, dessen Bücher immer kurze
Einführungen von rund 150--200 Seiten in ein schweizerisches Thema
darstellen, gedruckt in einem Format etwas kleiner als A5. Diese Themen
stammen oft aus der Geschichte, aber es gibt auch Biographien,
philosophische, politologische oder soziologische Beiträge. Oft liegt
der Fokus dabei noch spezifischer auf der französischsprachigen Schweiz,
der Romandie.

In diesem Kontext beschäftigt sich der Band \emph{La presse illustrée:
Une histoire romande} nun mit der Geschichte der illustrierten
Zeitschriften in der Romandie. In seiner Kürze kann das Buch nur einen
kleinen Teil der Zeitschriften intensiver besprechen (und bis auf eine
Ausnahme auch nur Titelblätter reproduzieren, obgleich die Besonderheit
dieser Zeitschriften die Doppelseite war, also die explizit auf
bestimmte Wirkungen hin geplante Anordnung von Bildern, Text und
graphischen Elementen auf zwei gegenüberliegenden Seiten). Zudem umfasst
der Band einen Zeitraum von 1842 -- mit der ersten Zeitschrift \emph{Le
cabinet de lecture} -- bis in die 1960er Jahre, was bei der Kürze des
Textes auch heisst, dass viele Themen, Entwicklungen und besprochene
Zeitschriften eher gestreift als eingehend behandelt werden. Es ist, dem
Anspruch von Savoir Suisse entsprechend, ein erster Überblick.

Als solcher ist der Band aber mit Gewinn zu lesen. Zu erfahren ist, dass
-- wenig überraschend -- die Entwicklung der illustrierten Zeitschriften
in der Romandie immer abhängig war von den sich entwickelnden
Möglichkeiten von Drucktechnik, Bildreproduktion und Photographie.
Zudem, dass immer Zeitschriften aus Frankreich und der Deutschschweiz
einen grossen Einfluss auf die Zeitschriften in der Romandie hatten:
Sowohl als Antrieb, spezifische Zeitschriften zu gründen und zu führen,
die weder einfach aus Frankreich, genauer Paris, noch aus der
Deutschschweiz, vulgo Zürich, geführt wurden, sondern die einen
expliziten \enquote{point du vue romand} einnahmen; als auch als
Konkurrenz, gegen die sich Zeitschriften an den Zeitschriftenkiosken in
Neuchâtel, Biel oder Delémont durchsetzen mussten. Gleichzeitig erfährt
man, dass es trotz dieser Herausforderungen in der Romandie immer wieder
eigene Zeitschriftenprojekte gab, teilweise alleinstehend, teilweise in
Zusammenarbeit mit deutschschweizer Verlagen (die mehrfach
französischsprachige Versionen von schon bestehenden deutschsprachigen
Zeitschriften publizierten, aber teilweise ergänzt durch explizite
Beiträge aus der Romandie). Wenn auch nicht als erste*r, so doch oft
kurz nach einer Entwicklung, die in Paris oder Zürich stattfand, zog
auch jemand in der Romandie nach (fast immer in Genf oder Lausanne,
manchmal aber auch in kleineren Städten wie Vevey).

Gleichzeitig zeigt der Band, wie insbesondere die beiden Weltkriege die
illustrierte Presse veränderten. Während des Ersten Weltkrieges
etablierte sich -- auch in der Romandie -- die Photoreportage, während
des Zweiten der explizite Einsatz von modern gestalteten illustrierten
Zeitschriften als Propagandainstrument (hier tat sich das
nationalsozialistische Deutschland mit der Zeitschrift \emph{Signal}
hervor, die in mehreren Sprachen für die neutralen Staaten und das
feindliche Ausland erschien; aber auch andere Ländern setzten diese Form
von Propaganda ein). Diese Entwicklungen wirkten sich dann auch auf die
Zeitschriften in den jeweils folgenden Friedenszeiten aus.

Interessant ist, dass Gianni Haver auch mehrfach darauf verweist, wie
schlecht der Überlieferungsstand gerade früher illustrierter
Zeitschriften in Bibliotheken der Romandie und Frankreichs ist. Von
vielen sind nur noch einige Ausgaben erhalten, bei anderen wissen wir
nicht, wann sie mit dem Erscheinen begannen oder endeten. (Dazu bei
trägt aber auch, dass es in der Schweiz und Frankreich keine der
deutschen Zeitschriftendatenbank vergleichbare Übersichten der in
Bibliotheken vorhandenen Zeitschriften gibt). Alles in allem bietet
dieser Band einen guten Überblick über die Entwicklung illustrierter
Zeitschriften im Allgemeinen und der Romandie im Besonderen. (ks)

\begin{center}\rule{0.5\linewidth}{0.5pt}\end{center}

Clavien, Alain (2025). \emph{L\textquotesingle argent de la presse
suisse: les années Publicitas, 1890--1990}. (Les routes de
l\textquotesingle histoire) Neuchâtel: Editions Livreo-Alphil
{[}gedruckt{]}

Werbung, beziehungsweise die Einnahmen aus ihr, war das gesamte 20.
Jahrhundert über die finanzielle Basis der Tages- und Wochenzeitungen
(in den marktwirtschaftlich organisierten Gesellschaften). Ihre
Bedeutung wuchs am Ende des 19. Jahrhunderts und ab den 1920er Jahren
stellten diese Einnahmen den Hauptteil des Budgets der meisten Zeitungen
dar. Organisiert wurde dies zumeist über Firmen, die zu Beginn
\enquote{Reklamebüros} hiessen und dann mit der Zeit ihre Bezeichnungen
veränderten. Diese organisierten die Werbung, welche in den Zeitungen
erschien, schlossen Verträge mit werbenden Firmen oder Werbeagenturen
und organisierten die damit zusammenhängenden Zahlungen. Die
Zeitungsverlage lagerten fast durchgehend diese wirtschaftliche Seite
ihrer Arbeit an Reklamebüros aus. In vielen Fällen \enquote{mieteten}
die Büros die Zeitungen auch über Jahrzehnte: Sie organisierten dann die
geschaltete Werbung exklusiv für eine Zeitung, zahlten dafür jährlich
garantierte Summen an die Verlage und erzielten ihren Profit aus der
Differenz zwischen diesen Zahlungen und den tatsächlichen Einnahmen.
Dafür kalkulierten die Zeitungen diese Summen in ihren Etat ein.

Allerdings: Obgleich die Arbeit dieser Reklamebüros für die Zeitungen
relevant waren, geschah sie oft im Verborgenen. Sie publizierten über
ihre tatsächliche Arbeit wenig und gewährten (beziehungsweise gewähren
auch heute noch) kaum Einblick darin. Allerdings liegen im Archiv des
Kantons Vaud / Waadt in Lausanne die Akten einer dieser Firmen, der
\enquote{Publicitas}, welche Ende des 19. Jahrhunderts begann, in der
Schweiz aktiv zu werden (1868 zuerst als Basler Zweigstelle einer
deutschen Firma, die dann 1890 zur eigenständigen Firma wurde und ab den
1930er Jahren in Lausanne ansässig war). Erst 2018 ging diese in den
Konkurs. Freigegeben sind diese, auch nicht vollständigen, Akten bislang
bis 1990. Sie geben einen Einblick in die Aktivitäten und auch die
Macht, welche diese Firma mit der Zeit auf den schweizerischen
Medienmarkt ausübte. In den 1930er Jahren wurde sie zur grössten dieser
Firmen, die dann auch noch mit den wenigen noch existierenden
Konkurrenten (andere hatte sie währenddessen aufgekauft) ein Kartell
bildete. Alain Clavien stellt die Geschichte dieser Firma in diesem
Buch, auf der Basis der überlieferten Akten, dar.

Dies ist erhellend: Einerseits zeigt sich, wie sehr die Presse in der
Schweiz über Jahrzehnte von den Werbeeinnahmen abhing. Sie waren nicht
etwa ein weiterer Einnahmeposten neben anderen, sondern deckten weit
über 50\,\% des Budgets der Zeitungsverlage ab. Wenn, dann waren die
Abonnements eine Nebeneinnahmequelle der Zeitungsverlage. Gleichzeitig
zeigt das Buch, wie sehr bestimmte Entwicklungen in den Medien von
politischen Entscheidungen dieser Firma abhingen. Beispielsweise wurden
sozialdemokratische Medien jahrzehntelang nicht bedient. Als diese in
den 1930er Jahren -- als die Sozialdemokratie sich deradikalisierte und
gleichzeitig in die schweizerische Parteienlandschaft eingefügt wurde --
von der Firma dann doch als \enquote{Kunden} zugelassen wurden, galt
dies weiterhin nicht für kommunistische oder andere radikale Blätter.
Noch in den 1970er Jahren verhinderte die Publicitas quasi, dass sich im
Kanton Valais / Wallis neben der dann bestehenden konservativen Zeitung
eine liberalere etablieren konnte, indem die Firma einen Kredit an die
Druckerei der letzteren vergab und dann, hinter der Bühne, bei der Firma
sanften Druck ausübte, die Arbeit für die liberale Zeitung einzustellen.
An sich etablierte sie sich über die Vergabe von solchen Krediten und
dem Ankauf von Firmenanteilen von Zeitungsverlagen selber praktisch als
Kontrollinstanz auf dem Zeitungsmarkt. Gleichzeitig stellte sich die
Firma nach aussen hin über Jahrzehnte als Stütze einer lokal
orientierten, lebendigen und offenen Zeitungslandschaft der Schweiz dar.

Die Firma prägte die Entwicklung der Zeitungslandschaft --
beispielsweise die Etablierung der \enquote{Kopftitel}, bei denen von
einem Verlag eine Anzahl von lokalen Titeln herausgegebene werden, die
grösstenteils die gleichen Beiträge enthalten und sich praktisch nur in
Titel, Lokalteil und Lokalanzeigen unterscheiden -- mit. Sie schaffte es
auch über eine lange Zeit, sich an Veränderungen, beispielsweise dem
Aufkommen des Fernsehens, anzupassen. Aber ab den 1990ern wurde es dann
offenbar immer schwieriger für sie, auf die gesellschaftlichen
Veränderungen bei der Mediennutzung und die Diversifikation von
Medientypen zu reagieren. Sie scheiterte dann, um es salopp
auszudrücken, praktisch am aufkommenden Internet.

An einigen Stellen ist das Buch recht langatmig, wenn es sehr
detailliert auf einzelne Geschäfte der Firma eingeht. Zudem ist es,
selbstverständlich, auf die Schweiz bezogen. Aber da es einen solchen
direkten Einblick in die Arbeit und Macht von Reklamebüros ansonsten
praktisch nicht gibt, ist es grundsätzlich ein Beitrag zur
Mediengeschichte, der sich auch auf andere Länder übertragen lässt. Es
ist wirklich zu empfehlen, insbesondere als Ergänzung zu
Mediengeschichten, die sich auf die Inhalte von Zeitungen fokussieren.
Hier geht es vor allem ums Geld, aber es zeigt sich dabei, wie wichtig
diese monetäre Seite immer war. (ks)

\begin{center}\rule{0.5\linewidth}{0.5pt}\end{center}

Matamoros, Isabelle (2025). \emph{Le pouvoir des lectrices: Une histoire
de la lecture au XIXe siècle}. Paris: CNRS Editions, 2025 {[}gedruckt{]}

Matamoros wertete für dieses Buch 64 Autobiographien, Tagebücher, Briefe
oder ähnliche Egodokumente von Frauen aus, die während des 19.
Jahrhunderts in Frankreich lebten. Dabei ging es darum, deren
Biographien als Leserinnen zu rekonstruieren: Wie haben sie lesen
gelernt, was haben sie gelesen, wie und wozu haben sie es gelesen? Es
ging darum, Gemeinsamkeiten herauszuarbeiten, welche ein spezifisches
Lesen von Frauen in dieser Zeit kennzeichneten. Eingebettet ist dies
alles in die Geschichte von Pädagogik, Erziehung und Unterricht, von
Literatur und Buchmarkt sowie zeitgenössischen Diskursen über das Lesen
(und seine Gefahren). Sicherlich: Frauen, die solche Egodokumente
verfassen, waren selber eine kleine Gruppe, obgleich es im 19.
Jahrhundert normaler war, als vielleicht heute. Matamoros versucht aber,
Frauen aus möglichst verschiedenen Schichten, Milieus und
Lebensgeschichten einzubeziehen. Dabei hilft, dass in der
Nationalbibliothek in Paris eine eigene Abteilung existiert, die aktiv
Tagebücher sammelt.

Die Lesebiographien zeichnen sich dann auch durch eine hohe Varietät
aus. Frauen lernten auf sehr verschiedenen Wegen zu Lesen, oft durch
ihre Familien selber, da eine geregelte Schulbildung sich erst im 19.
Jahrhundert durchsetzte. Gleichzeitig zeigte sich, dass pädagogische
Handbücher, die für den Unterricht daheim geschrieben wurden, einen
Einfluss auf Familien in allen sozialen Schichten hatten. Fast alle
Frauen waren dann die ganze Zeit des Lebens begleitet von Autoritäten,
welche ihr Lesen kontrollieren wollten, sowohl ihren Familien und
Ehepartnern als auch von Staat, Kirche und Öffentlichkeit. Das bezog
sich darauf, was sie lasen und wie sie lasen. So wurden einigen Frauen
angehalten, bestimmte Bücher nicht zu lesen, während andere angehalten
wurden, Bücher in einer ganz spezifischen Form (Matamoros nennt es
\enquote{Arbeit}, also nach bestimmten Listen, Vorgaben und so weiter)
zu lesen. Aber gleichzeitig bestimmten Frauen ihre Lesebiographie
selber. Sie lasen praktisch \enquote{um die Vorgaben herum}, sie lasen,
was sie nicht lesen sollten und sie wählten selber aus. Interessant ist,
wie instrumentell viele Frauen Bücher benutzten. Sie bereiteten sich zum
Beispiel auf Museumsbesuche oder Reisen aktiv durch Lektüre vor, sie
bildeten sich selber aktiv mit Büchern weiter und sie nutzen Literatur,
um die Welt um sie herum zu verstehen. Dabei waren sie vom Medienangebot
-- also den Büchern und Magazinen, die verkauft oder verliehen wurden
sowie den Zeitungen, die während des Jahrhunderts immer bedeutender
wurden -- beeinflusst, aber wählten aus diesem Angebot selber aus.

Das Buch ist offensichtlich für Menschen geschrieben, die ein
Grundwissen über die Geschichte Frankreichs im 19. Jahrhundert haben.
Die politischen und gesellschaftlichen Ereignisse werden immer nur kurz
erwähnt, aber nicht beschrieben. Es wird also zum Beispiel
vorausgesetzt, dass die Abfolgen von Monarchie und Republik,
Revolutionen und Restaurationen bekannt sind. Fokus sind immer die
Lesebiographie von Frauen, also die Alltagsgeschichte, die sich vor
diesem Hintergrund abspielte. (ks)

\begin{center}\rule{0.5\linewidth}{0.5pt}\end{center}

Antonutti, Isabelle (2024). \emph{Bâtisseuses de la lecture publique:
Une histoire des premières bibliothéquaires, 1900-1950}. Villeurbanne:
Presse de l'Enssib, 2024,
\url{https://doi.org/10.4000/books.pressesenssib.18637}

Die Geschichte des ersten Bibliothekarinnen in den Öffentlichen
Bibliotheken Frankreichs -- das Thema dieses Buches -- ist nur
bruchstückhaft dokumentiert und nachvollziehbar. Es gibt einige Hinweise
auf Biographien, insbesondere da in Frankreich früh eine Schule (heute
die Enssib, in deren Verlag dieses Buch erschien) für Bibliothekar*innen
eingerichtet wurde, deren Archiv sich erhalten hat. Aber ansonsten sind
es Bruchstücke, die sich finden lassen: Einige Beiträge in
Zeitschriften, Erwähnungen in Jahresberichten, Bibliothekarinnen auf
Photos.

Antonutti hat diese Bruchstücke zusammengetragen und versucht, eine rote
Linie zu finden. So ist das Buch auch gestaltet. Es gibt zwei Teile: Im
zweiten finden sich Biographien von Bibliothekarinnen aus dem frühen 20.
Jahrhundert, soweit sie rekonstruiert werden konnten. Im ersten Teil
findet sich ein Text, welcher die rote Linie beschreiben will, aber
immer wieder unterbrochen wird durch einzelne Dokumente oder Sammlungen
von Zitaten -- explizit graphisch abgesetzt -- von und über die
Bibliothekarinnen. Auch dieser Text selber wechselt ständig: An Stellen
macht er übergreifende Aussagen, an anderen thematisiert er, wie wenig
konkrete Quellen es gibt (und das sie kritisch zu werten sind) und an
vielen anderen Stellen beschreibt er einzelne Dinge, die dokumentiert
sind, ganz genau. Damit vermittelt der Text selber den Eindruck einer
Dokumentensammlung, die noch recht unvollständig und unausgeglichen ist.

Grundsätzlich wird klar, dass Bibliothekar*innen in den Öffentlichen
Bibliotheken in Frankreich am Anfang des 20. Jahrhunderts
\enquote{auftauchen} und dass dies im europäischen Rahmen eine normale
Entwicklung war. Es scheint, dass einzig in den USA einige Jahrzehnte
zuvor die \enquote{männliche Domäne} Bibliothek auch für Frauen geöffnet
worden war. (ks)

\begin{center}\rule{0.5\linewidth}{0.5pt}\end{center}

Höppner, Stefan (2022). \emph{Goethes Bibliothek: Eine Sammlung und ihre
Geschichte}. (Zeitschrift für Bibliothekswesen und Bibliographie,
Sonderbände, 125) Frankfurt am Main: Vittorio Klostermann, 2022
{[}gedruckt{]}

Die Bibliothek Johann Wolfgangs von Goethe ist grösstenteils in Weimar
in seinem langjährigen Arbeits- und Lebensort im \enquote{Haus am
Frauenplan} erhalten. Goethe war einer der ersten Autor*innen, welcher
aktiv den eigenen Nachlass und Nachruhm zu planen begann. Damit hatte er
auch Erfolg, wie die Verehrung als einer der \enquote{deutschen
Nationaldichter} in den auf seinen Tod folgenden Jahrhunderten zeigte.
Dies motivierte auch ständig Forschungen zum Leben und Werk Goethes, die
auch auf die konkrete Bibliothek zugriffen (oder zugreifen wollten).
Heute ist die Bibliothek -- allerdings in einer renovierten, durch ein
Gitter ergänzt Variante und nicht mehr so aufgestellt, wie sie es zum
Tod Goethes war, zudem mit teilweise fehlenden, teilweise in
Einzelstücken auch erweitertem Bestand -- museal besuchbar und auch ein
beliebtes Bild für Postkarten. Was die Bibliothek auszeichnet, ist, dass
sie als Arbeitsbibliothek genutzt wurde und gerade keine
Repräsentationszwecke hatte. Zudem wurde die Bibliothek mehrfach
katalogisiert, nicht immer erfolgreich und vollständig, zuletzt im
Projekt \enquote{Goethe Digital} (2014--2024,
\url{https://vfr.mww-forschung.de/web/goethedigital/start}), in dessem
Rahmen zudem auch Visualisierungen, beispielsweise zur Nutzung anderer
Bibliotheken durch Goethe oder zu Herkunft der Bücher in Goethes
Bibliothek, erstellt, Ausstellungen organisiert, Tagungen durchgeführt
und Publikationen veröffentlicht wurden.

Dieses Buch stellt eine Art Rundumschlag zur Bibliothek und zum Projekt
dar. Stefan Höppner leitete dabei das genannte Projekt. Gegliedert ist
das Buch in zwei Teile, einen historischen und einen systematischen. Im
ersten Teil wird nicht nur die Geschichte der Bibliothek, sondern einer
ganzen Anzahl von weiteren Themen, die mit der Bibliothek zu tun hatten,
geschildert. Es geht beispielsweise um die Goethe-Rezeption im deutschen
Kaiserreich, der Weimarer Republik, im Nationalsozialismus und der DDR
und darum, wie über die Jahrhunderte immer wieder neue Kataloge
versprochen, aber dann doch nicht abgeschlossen wurden. Gleichzeitig
schildert der Text die allgemeine Entwicklung von Bibliotheken oder auch
Auseinandersetzungen zwischen der Familie Goethe und den verschiedenen
deutschen Staaten. Im systematischen Teil werden vor allem
Forschungsergebnisse aus dem Projekt berichtet, beispielsweise über
verschiedene Exlibris, über die Frage, was eine Sammlung ist und ob die
Bibliothek Goethes auch als Sammlung verstanden werden kann oder wer
Goethe Bücher schickte. Im zweiten Teil werden auch Visualisierungen und
Tabellen aus diesen Forschungen verwendet.

Das ist alles thematisch viel, unterschiedlich aufbereitet und
beantwortet sehr verschiedene Fragen. Und das ist das Problem des
Buches: Es ist einfach zu viel, was hier behandelt wird. Problematisch
ist dabei wohl vor allem, dass es als ein zusammenhängender Text
präsentiert wird. Dadurch gehen aber die einzelnen Themen und Fragen
ständig unter. Es vermittelt den Eindruck, als hätte hier alles, was im
Projekt irgendwie erarbeitet wurde, in ein Buch gepackt werden sollen.
Hat man am Anfang beispielsweise über die Beziehung von Goethe zum
Fürstenhaus in Weimar gelernt, kommt man zu Erbstreitigkeiten oder zur
Beziehung des nationalsozialistischen Regimes zur Bibliothek. Gelangt
man dann beim systematischen Teil und der Darstellung an, welche Bücher
woher stammen, ist es schwer, noch diese später gegebenen Informationen
-- die ja eigentlich mehr über die eigentliche Bibliothek sagt --
wertzuschätzen. Es hätte dem Buch gut getan, entweder auf mehrere
Publikationen verteilt zu werden oder aber als Sammlung eigenständiger,
abgeschlossener Beiträge präsentiert zu werden. In der publizierten
Variante ist es aber überwältigend. Es ist alles irgendwie relevant und
interessant, aber es sollte immer nur in kleinen Abschnitten gelesen
werden. (ks)

\begin{center}\rule{0.5\linewidth}{0.5pt}\end{center}

Rausch, Kerstin Annika (2024). \emph{Antike Bibliotheken: Bildungsorte
zwischen Symposion, Religion und gesellschaftlicher Repräsentation}.
(Göttinger Studien zur Mediterranen Archäologie, 15) Rahden / Westfalen:
Verlag Marie Leidorf, 2024 {[}gedruckt{]}

Diese Dissertation sammelte zusammen, was aus antiken Quellen, deren
Interpretation in den letzten Jahrhunderten, und archäologischen Funden
über Bibliotheken in der griechischen und römischen Antike bekannt ist.
Das ist alles ein Überblick, recht leidenschaftslos, aber dafür
umfassend und möglichst vollständig berichtet. Als solcher präsentiert
er keine grossen Überraschungen. Rausch entdeckt keine neuen Funde und
stellt auch keine neuen Thesen auf, sondern zeigt, dass die Bibliotheken
unterschiedliche Funktionen hatten (Bildung, Teil von
Freizeitgestaltung, Repräsentation) und deshalb auch in verschiedene
Kontexte eingebunden waren. Dieser Eindruck eines Überblicks wird durch
das Format des Buches (A4, relativ dickes Papier, harte Bindung, recht
einfacher Satz) unterstützt. Es sieht aus, wie ein funktionales
Schulbuch.

Die interessanteste Information findet sich gleich auf der ersten Seite
des Vorwortes: Rausch hat im Rahmen dieser Arbeit eine, wohl möglichst
vollständige, Datenbank zu antiken Bibliotheken aufgebaut, die am
eScience Center Tübingen verwaltet wird. Allerdings ist diese Datenbank,
die Grundlage dieser Arbeit war und offensichtlich eine gute Grundlage
weiterer Forschungen sein kann, nicht einfach zugänglich. Dafür muss
erst Kontakt zum Archäologischen Institut der Georg-August-Universität
Göttingen aufgenommen werden: Eine für heutige Verhältnisse erstaunliche
Regelung. (ks)

\begin{center}\rule{0.5\linewidth}{0.5pt}\end{center}

Rathkolb, Oliver (2025). \emph{Kontrollierte Freiheit: Die Alliierten in
Wien -- Kulturpolitik 1945--1955}. Wien: Wien Museum, Residenz Verlag,
2025. {[}gedruckt{]}

Nach dem Ende des Zweiten Weltkrieges wurden auch Österreich und dann
noch einmal die Stadt Wien in Besatzungszonen unterteilt, die von den
Siegermächten kontrolliert wurden. Anders als in Deutschland endete
diese Besatzungszeit -- obgleich der Kalte Krieg schon begonnen hatte --
1955 mit dem Staatsvertrag, mit welchem die Souveränität Österreichs
wiederhergestellt wurde. Das Land band sich dann explizit
\enquote{Richtung West}. Aus Anlass der Wiederkehr dieses Jubiläums fand
in Wien 2025 eine Ausstellung über die Kulturpolitik der Siegermächte
und deren zeitgenössischer Wirkung statt. Der vorliegende Band ist die
Begleitpublikation dazu.

In grossem Format (A4), lockerem Layout, aber auch mit grossem Font und
vielen Abbildungen, stellen Beiträge von je rund zehn Seiten einzelne
Aspekte dieser Kulturpolitik wie Magazine, Zeitungen und Radiostationen
der Alliierten, Ausstellungen, Förderung von Theater und Kino sowie
anderes dar. Teilweise mit Fokus auf eine Besatzungsmacht, teilweise in
Rundschau über alle vier. Zu lernen ist dabei unter anderem, dass es
neben der Entnazifizierung ein gemeinsames Ziel aller Besatzungsmächte
war, eine österreichische Identität aufzubauen, die sich explizit von
der deutschen beziehungsweisen deutschnationalen (die Österreich als
Teil Deutschland ansah) unterschied. Insbesondere Frankreich setzte
darauf, eine Traditionslinie der Verbindung zwischen Österreich und
Frankreich zu betonen. Ausserdem wird klar, dass die USA den Kampf um
die Köpfe der Massen gewann und insbesondere die Sowjetunion ihn verlor.

Erwähnt wird dieser Band hier aber wegen einer Leerstelle, die
bezeichnend ist für solche Geschichtsschreibungen: Es gibt zwar einen
eigenen Beitrag über Bibliotheken in Wien in dieser Zeit, der auch
grundsätzlich interessant ist. (Markus Stumpf. \emph{\enquote{Bücher
gehören zweifellos nicht zu den Möbeln.}: Kampfzone Bibliothek --
Reorientierung durch Entnazifizierung und Büchergeschenke}, ebenda:
158--169) Dieser beschäftigt sich vor allem damit, wie die grossen
Wissenschaftlichen Bibliotheken (zu denen auch die Stadtbibliothek Wien
zählte, die damals noch nicht zum Netzwerk der städtischen Büchereien
gehörte) von nationalsozialistischer Literatur gesäubert wurden
(zumindest teilweise) und wie die Bestände wieder aufgebaut wurden.
Zudem gibt es einen Artikel explizit zur Verlagspolitik (Günter Stocker:
\emph{Die Alliierten und die österreichische Literatur}, ebenda:
170--179). Aber was nicht thematisiert wird, sind die Bibliotheken,
welche eine grosse Masse der Wiener*innen erreichte, also die
städtischen Büchereien und die offenbar zahlreichen Leihbüchereien und
Lesesäle, die von den Alliierten eingerichtet wurden. Das ist
einigermassen frustrierend: Immer wieder wird in anderen Beiträgen
erwähnt, dass eine der ersten Aktivitäten der Siegermächte war, eigene
Leseräume und Leihbüchereien einzurichten, oft noch bevor andere
Einrichtungen eröffnet oder neue Medien herausgegeben wurden. Teilweise
passierte das im jeweils von ihnen kontrollierten Stadtgebiet, teilweise
in den Kultur- und Informationszentren, welche von den Siegermächten
aufgebaut wurden. Ganz offensichtlich fanden sie diese Aktivität damals
wichtig und potentiell einflussreich. Zudem unterschieden sie offenbar
ständig die beiden Einrichtungen Leseraum und Leihbücherei. Aber: Was in
diesen Einrichtungen passierte, welche Medien in ihnen standen, wie sie
genutzt wurden und von wem, warum die beiden Einrichtungen bei den
Alliierten so explizit getrennt wurden -- das scheint im ganzen Band
nicht zu interessieren. Offenbar sind \enquote{grosse Themen} wie die
Ausstellungen damals moderner Kunst oder auch -- um eine im Band
besprochene Massenkultur zu nennen -- der Sport für die
Geschichtsschreibung relevanter als die wohl breit wirksamen, aber
kleinen Büchereien für die breite Masse. (Was, wie in dieser Kolumne
hier schon mehrfach festgestellt wurde, nicht nur für diesen Band gilt,
sondern offenbar für viele Geschichtsschreibungen.) Das ist eine
ärgerliche blinde Stelle, da es thematisch in diese Ausstellung und
diesen Band sehr gut gepasst hätte. (ks)

\section{4. Weitere wissenschaftliche Medien (Konferenzberichte,
Abschlussarbeiten)}\label{weitere-wissenschaftliche-medien-konferenzberichte-abschlussarbeiten}

{[}diesmal keine Beiträge{]}

\section{5. Populäre Medien (Social Media, Zeitungen, Radio,
TV)}\label{populuxe4re-medien-social-media-zeitungen-radio-tv}

Fischer, Georg (2026). „\emph{Die Daten fließen nur in eine Richtung,
nämlich von der Wissenschaft zu Elsevier'' -- Yuliya Fadeeva zum
Datentracking durch Verlage}.
\url{https://irights.info/artikel/datentracking-elsevier-yuliya-fadeeva-interview/32777}
{[}abgerufen 06.03.2026{]}.

Im Interview wird ein kurzer Überblick über die Geschäftspraktiken von
Elsevier und anderen Großverlagen -- oder Ex-Verlagen, wie sie von der
Interviewpartnerin bezeichnet werden -- gegeben. Schwerpunkte sind dabei
der Wandel vom Verlag hin zum Data-Analytics-Unternehmen, die
Ungleichheiten und Abhängigkeiten im Verhältnis von Forschung und
Großkonzernen sowie mögliche Wege aus dieser Schieflage. (eb)

\begin{center}\rule{0.5\linewidth}{0.5pt}\end{center}

Forester, Brett ; Deer, Ka'nhehsí:io ; Luke, Marnie ; Seglins, Dave
(2026). \emph{How RCMP spies infiltrated the 1970s Indigenous rights
movement}. CBC Investigates,
\url{https://www.cbc.ca/news/indigenous/rcmp-spies-1970s-indigenous-rights-9.7134112}
(24.03.2026)

Seglins, Dave (2026). \emph{Getting access to RCMP
\textquotesingle Native extremism\textquotesingle{} files took 4-year
fight}. CBC Investigates,
\url{https://www.cbc.ca/news/indigenous/rcmp-access-1970s-spying-files-9.7143553}
(27.03.2026)

Forester, Brett (2026). \emph{The RCMP, the spy and the betrayal of
national chief George Manuel}. CBC Investiges,
\url{https://www.cbc.ca/news/indigenous/rcmp-security-infiltration-george-manuel-9.7145083}
(28.03.2026)

St-Pierre, Sarah (2026). \emph{Apology \textquotesingle not
enough\textquotesingle: Dene leaders want to see the RCMP spying files.}
CBC North,
\url{https://www.cbc.ca/news/canada/north/apology-not-enough-dene-leaders-want-see-rcmp-spying-files-9.7145344}
(30.03.2026)

Die kanadische Polizei hat in den 1970er Jahren massiv aktivistische
Gruppen von First Nations ausspioniert und unterwandert. Diese Gruppen
setzten sich damals (und setzen sich teilweise auch heute noch) für die
Rechte der First Nations ein, stellten Forderungen auf, wie die, dass
First Nations gleichberechtigt mit der kanadischen Regierung über ihre
Zukunft verhandeln und eigene Entscheidungen über ihre Entwicklung
treffen können sollten. Dies alles geschah in einer bewegten Zeit, die
geprägt war von Bürgerrechtsbewegungen, auch in Kanada, und von einem
steigenden Aktivismus von First Nations oder, wie sie sich in den USA
oft nennen, American Natives. Teilweise lief dies unter dem Schlagwort
\enquote{Red Power}. Sowohl in den USA als auch in Kanada ging dies mit
teilweise aufsehenerregenden Besetzungen (unter anderem von Alcatraz
oder Wounded Knee) und radikalen Protesten einher. In gewisser Weise
spiegelte das Vorgehen der kanadischen Polizei gegen die Aktivist*innen
der First Nations das Cointelpro-Programm, mit dem das FBI in den USA
die afroamerikanische Bürgerrechtsbewegung, aber auch das
\enquote{American Indian Movement}, unterwanderte.

Offenbar hatten viele Aktivist*innen damals schon die Vermutung, dass
sie überwacht wurden. Aber erst 2026 deckte der öffentliche Rundfunk CBC
(Canadian Broadcasting Cooperation) den Umfang und die Breite dieser
Überwachung auf. In einer wachsenden Anzahl von Artikeln dokumentierte
der CBC seit Ende März 2026 die Arbeit der Polizei, interviewte
Betroffene und berichtete auch über Reaktionen von offizieller Seite,
bis hin zum Premierminister, der eine offizielle Entschuldigung und
Aufklärung ankündigte.

Die Hauptquellen für diese Berichte stellen Akten dar, die im
kanadischen Nationalarchiv (Library and Archives Canada) liegen, aber --
wie Dave Seglins in einem eigenen Artikel berichtete -- erst durch
mehrere Anfragen mit Verweisen auf das kanadische
Informationsfreiheitsgesetz und Beschwerden an die betreffende
Ombudsstelle zugänglich gemacht wurden. Das Archiv hielt die Akten
zurück und verwies, weil sie die nationale Sicherheit betreffen würden,
an den kanadischen Geheimdienst. Dieser verwies wieder auf das Archiv
zurück, das dann später dem CBC-Team die betreffenden Akten zugänglich
machte. Im Laufe der Veröffentlichungen äusserte es dann auch Interesse,
einen noch breiteren Zugang zu gewähren. Ein Grund, der vom Archiv gegen
die Veröffentlichung der Akten vorgebracht wurde, war der, dass aus den
Akten offenbar auch ersichtlich sein kann, welche Personen konkret die
Spionage durchführte, also wer Berichte lieferte oder gar aktiv Gruppen
unterwanderte.

Mehrere Interviewte, die (oder deren heute verstorbenen Eltern) damals
von der Überwachung betroffen waren, äusserten dagegen explizit die
Forderung, dass die Akten veröffentlicht werden sollten. Damit alleine
sei es nicht getan, genauso wenig wie mit einer Entschuldigung von
staatlicher Seite. Aber es wäre ein Anfang zur Aufarbeitung. Ein wenig
erinnert diese Situation an ähnliche Positionen, wie sie nach dem Ende
der DDR im Bezug auf Stasi-Akten vertreten wurde, obgleich diese
Parallele -- aus der man vielleicht etwas über den möglichen Umgang mit
solchen Akten lernen könnte -- (bislang) in Kanada nicht erwähnt wurde.
Gleichzeitig erinnert dieses Beispiel aber auch daran, welche
Verantwortung Archive tragen können. (ks)

\begin{center}\rule{0.5\linewidth}{0.5pt}\end{center}

Rüegger, Samuel (2026). \emph{LibreOffice am Abgrund: Wie die Document
Foundation ihre eigenen Gründer vor die Tür setzt}. Linux Guides
Community,
\url{https://forum.linuxguides.de/core/index.php?article/54-libreoffice-am-abgrund-wie-die-document-foundation-ihre-eigenen-gr\%C3\%BCnder-vor-die/}
{[}abgerufen 09.04.2026{]}.

Der Blogpost gibt einen Einblick in die Schwierigkeiten, die mit der
Betreuung einer großen und erfolgreichen Open-Source-Software
einhergehen, wobei Themen der Finanzierung ebenso angerissen werden, wie
Problematiken des Steuerrechts und ethische Grundlagen der
Open-Source-Bewegung. Exemplarisch festgemacht wird es an den 2020
begonnenen Geschehnissen rund um LibreOffice. Die gemeinnützige, in
Berlin ansässige Document Foundation, die über die Markenrechte an
LibreOffice verfügt, hat sich heftig mit maßgeblichen Entwicklern der
Software verkracht -- bis hin zu Rechtsstreitigkeiten -- was für die
Zukunft der dieser freien Microsoft-Office-Alternative nichts Gutes
befürchten lässt. (eb)

\section{6. Weitere Medien}\label{weitere-medien}

Glättli, Samuel (2024). \emph{Globine baut eine Bibliothek}. Zürich:
Globi Verlag ; Orell Füssli Verlag, 2024 {[}gedruckt{]}

Globine und Globi sind zwei Comicfiguren, welche zur Alltagskultur der
deutschsprachigen Schweiz und Liechtensteins gehören, aber ausserhalb
dieser Gebiete kaum bekannt sind. (Vergleichbar vielleicht mit den
Abrafaxen in Ostdeutschland.) Die Comics gehören seit Jahrzehnten zur
Kindheit in der Schweiz und Liechtenstein, die Bücher stehen hier wohl
in den meisten Kinderzimmern, sind aber auch in praktisch allen
Gemeindebibliotheken oder Buchhandlungen zu finden.

Ursprünglich wurde Globi 1935 als Werbefigur für die schweizerische
Kaufhauskette Globus geschaffen, aber diese Verbindung ist schon lange
nicht mehr gegeben. Globine wurden 1988 eingeführt. In ihren Geschichten
fungieren beide immer wieder als Figuren, welche oft eine Institution
oder Sehenswürdigkeit in der Schweiz vorstellen (die Rettungsflugwacht,
den Zoo, das Kuhrennen oder -- in seiner rätoromanischen Variante eines
der prominentesten Bücher in dieser Sprache -- den Nationalpark).
Grundsätzlich sind beide immer unternehmenslustig, aktiv und positiv
gestimmt. In den Geschichten lösen sich alle Probleme durch gemeinsamen
Einsatz, beispielsweise der Tiere im Nationalpark oder aller
Bewohner*innen eines Dorfs, mit einem positiven Ergebnis auf. Dabei
werden verständlicherweise viele Herausforderungen und deren Lösung
vereinfacht und ihre Komplexität reduziert. Die Darstellung ist jeweils
wenig komplex, die Bildtafeln sind, der Zielgruppe Kinder gerecht,
jeweils bunt und relativ einfach aufgebaut.

Das ist relevant für den Kontext des Bandes \enquote{Globine baut eine
Bibliothek}, welcher 2024 erschien. Er hebt sich in Stil und Geschichte
nicht von den anderen Globine / Globi-Bänden ab, aber er ist sofort --
erwartungsgemäss -- zu einem Hit in den Bibliotheken der Deutschschweiz
und Liechtensteins geworden. In praktisch allen steht er im Bestand. Der
Rezensent hat ihn in den letzten Monaten zudem mehrfach in den Regalen
verschiedener Bibliotheken (inklusive Wissenschaftlicher) hervorgehoben
gesehen. Er ist aber auch oft entliehen (der Rezensent griff extra auf
den Bestand der Stiftsbibliothek St.~Gallen zurück, um es ohne das
schlechte Gewissen, dass er es wohl einem Kind fortnimmt, für einige
Tage auszuleihen). Es ist in kurzer Zeit zu einem Klassiker geworden --
einer, der ausserhalb dieser beiden Länder wohl unbekannt ist.

Das Globine eine Bibliothek baut (mit Hilfe von tierischen Freund*innen
wie einem Eichhörnchen und einer Eule sowie offenbar der gesamten
Bevölkerung ihres Dorfes Schnurzikon und der Bürgermeisterin), ist
motiviert davon, dass daheim ihr eigenes Buchregal auseinanderfällt, sie
beim Besorgen neuen Holzes (sie baut grundsätzlich alles selber) Kinder
trifft, die mit dem Rad die Bibliothek in der Nachbargemeinde besuchten,
sich aber über die weite Strecke beschweren, und das zudem in der
Dorfmitte von Schnurzikon ein leeres Gebäude steht. Dies alles bringt
Globine zu der Idee, dass man dort eine Bibliothek einrichten müsste.
Nachdem diese zuerst mit grossem Erfolg startet, aber nur mit Büchern
und händischer Ausleihe ausgestattet ist, wird sie von Globine
modernisiert, die dabei auch Ideen anderer Figuren aufgreift. Sie wird,
wie das im bibliothekarischen Diskurs heissen würde, zu einem
\enquote{Zentrum des Dorfes} mit Podcaststudio, Werkstatt,
Kulturprogramm (Vorleseveranstaltungen und Konzert), Buchausstellung,
gemütlichen Ecken und Kaffeeautomat. Gleichzeitig wird sie technisch
aufgerüstet. Globine baut erst eine Selbstausleihe, dann einen Roboter,
welcher die Einstellarbeiten übernimmt. Die Geschichte endet damit, dass
die Bibliothek nach einer Brandkatastrophe (und dem Finden von
antiquarischen Büchern, die zum Teil veräussert werden, um die
notwendigen Mittel für den Neubau aufzubringen) wieder eingerichtet und
vom ganzen Dorf genutzt wird.

Selbstverständlich: Die Geschichte übergeht viele Dinge, die realen
Bibliotheken Kopfzerbrechen machen. Beispielsweise wird die Bibliothek
zuerst mit Buchspenden aufgebaut, die in der Realität mehr Problem als
Lösung sind. Es wird auch -- obwohl die Bürgermeisterin direkt die
Bibliothek unterstützt -- nicht geklärt, wer überhaupt den Betrieb
finanziert. Das alles ist dem Genre geschuldet. Aber sonst zeichnet das
Buch das Bild einer Bibliothek, in der Bücher und das Lesen im
Mittelpunkt stehen, welche aber sonst auch eine grosse Breite an
Angeboten hat; die einen Ort darstellt, in der viele Menschen (und
Tiere) konfliktfrei ihren jeweils eigenen Interessen nachgehen können
und die zudem modern ist.

Der Band ist in sich abgeschlossen, verweist aber auch auf andere Bände
der Reihe zurück. Zudem kehren aus diesen zahlreiche Figuren und
Personen wieder. Eine Kenntnis früherer Bände fördert also das
Verständnis. Zudem gibt es in der Geschichte eine Reihe von Verweisen
auf den schweizerisch-liechtensteinischen Alltag, die wohl für
Leser*innen aus anderen Ländern nicht wirklich verständlich sind.
Trotzdem eignet sich der Band auch für die Kinderabteilungen der
Bibliotheken anderer Länder. (ks)

%autor

\end{document}